\NeedsTeXFormat{LaTeX2e}[1996/06/01]
\documentclass[11pt,twoside,makeidx,hidelinks,]{memoir}
\usepackage{layouts}[2001/04/29]
\def\isfinal{1}

% Packages
%=========
\usepackage{verbatim}
\usepackage[colorlinks=false]{hyperref}
% \usepackage[utf8x]{inputenc}
\usepackage{listings}
\usepackage{
  fancyvrb,
  graphics,
  url,
  rotating,
  lipsum,
  microtype,
  fancybox,
  glossaries,
  titling,
  xspace,
}
\renewcommand*{\thefootnote}{\fnsymbol{footnote}}

% Boolean-ness
%=============
\newcommand{\onlyfinal}[1]{\ifdefined\isfinal{#1}\fi}%
\newcommand{\notfinal}[1]{\ifdefined\isfinal\else{#1}\fi}%
\newcommand{\whenfinal}[2]{%
  \ifdefined\isfinal%
    #1%
    \else%
    #2%
  \fi}%

%
% Page Layout
%============


\setstocksize{9in}{6in}
\setulmarginsandblock{0.8in}{0.5in}{*}
\settrims{0in}{0in}
\setlrmarginsandblock{0.8in}{0.5in}{*}
\settrimmedsize{9in}{6in}{*}
\settypeblocksize{7.7in}{4.6in}{*}

\setmarginnotes{0.1pt}{0.2in}{\onelineskip}
\setheadfoot{\onelineskip}{2\onelineskip}
\setheaderspaces{*}{2\onelineskip}{*}
\linespread{1.15}

\checkandfixthelayout

\setlength{\headsep}{8pt}

\newlength{\AfterFancyBreakMin}
\setlength{\AfterFancyBreakMin}{3\baselineskip}
\setlength{\parindent}{1.5 em}
\makeatletter
\g@addto@macro\quote{\linespread{1}\small{}}
\g@addto@macro\quotation{\linespread{1}\small{}}
\g@addto@macro\caption{\linespread{1}\footnotesize{}}
\makeatother
\tolerance 1414
\hbadness 1414
\setlength{\emergencystretch}{3em}
\righthyphenmin 62
\emergencystretch 1.5em
\hfuzz 0.3pt
\widowpenalty=10000
\clubpenalty=10000
\setlength{\parindent}{1.5 em}

% Page Heading
%-------------
\makepagestyle{ebook}
\makeoddhead{ebook}{}{\textsf\theauthor}{\normalfont\thepage}
\makeevenhead{ebook}{\normalfont\thepage}{\textsf \thetitle}{}

\aliaspagestyle{title}{empty}

% Text Tweaks
%-------------
\def\thinspace{\kern .16667em }
\newcommand{\ellip}{\ldots}
\newcommand{\Dash}{\thinspace---\thinspace}
\newcommand{\dash}{\tk{--Wrong Dash--}}

% Fonts
%------
\usepackage{fontspec}
\usepackage{xunicode}
\usepackage{xltxtra}
\defaultfontfeatures{Mapping=tex-text}
\setmainfont{Georgia}
\setsansfont{Tahoma}
\setmonofont{Courier}
\newfontinstance\scshape[Letters=SmallCaps, Numbers=Uppercase]{Courier}
\DeclareOldFontCommand{\tt}{\normalfont\ttfamily}{\mathtt}
\ifdefined\disableligature
  \DisableLigatures[f]{encoding = *, family = * }
\fi

% Dates
%-------
\usepackage{datetime}
\newdateformat{yearonly}{\THEYEAR}
\newdateformat{usmonthyear}{\monthname[\THEMONTH], \THEYEAR}
\newcommand{\monthyear}{%
  \ifcase\month\or January\or February\or March\or April\or May\or June\or
  July\or August\or September\or October\or November\or
  December\fi\space\number\year
}

% Document Headings
%==================

% Chapter
%-----------
\setsecheadstyle{\bfseries\large\centering}
\newcounter{renewsection}
\newcommand{\Chapter}[1]{
  \setcounter{renewsection}{\value{section}}
  
  \chapter{}
  
  \setcounter{section}{\value{renewsection}}
}
\chapterstyle{chappell}


% Section
%---------
\renewcommand{\thesection}{\arabic{section}.}

% Novel Tweaks
%=============

\newenvironment{xdedication}%
        {\vspace*{6ex}\begin{quotation}\begin{center}\begin{em}}%
        {\par\end{em}\end{center}\end{quotation}}%

\newenvironment{dedication}
  {\clearpage           % we want a new page
   \thispagestyle{empty}% no header and footer
   \vspace*{\stretch{1}}% some space at the top
   \itshape             % the text is in italics
   \raggedleft          % flush to the right margin
  }
  {\par % end the paragraph
   \vspace{\stretch{3}} % space at bottom is three times that at the top
   \clearpage           % finish off the page
  }

% Scene Formatting
%-----------------

\newcommand*\thoughtbreak{
  \begin{center}
    \vspace{1cm}
  \noindent\rule{0.5\textwidth}{0.4pt}
    \vspace{1cm}
  \end{center}
}

\newcommand{\thought}[1]{\emph{#1}}
\newcommand{\email}[1]{\texttt{#1}}
\newcommand{\status}[1]{\notfinal{\textbf{Status:} #1\newline}}
\newcommand{\synopsis}[1]{\notfinal{\line(1,0){320}\newline\textbf{What's going on here?} #1\newline
\line(1,0){320}\newline}}
\newcommand{\slug}[1]{\notfinal{\MakeUppercase{#1}\newline\newline}}
\newcommand{\slugline}[2]{\notfinal{\section{#1}}\onlyfinal{\section{}}\slug{#2}}
\newcommand{\interlude}[1]{\plainfancybreak{\AfterFancyBreakMin}{1}{#1}}

% \newcommand{\scene}[3]{%
%   \notfinal{\section{#1}\slug{#2}}
%   \input{#3}
%   \pbreak
% }%

\newcommand{\smee}[1]{\${-} \textsc{\small#1}}
\whenfinal%
  {\newcommand{\pbreak}{\pfbreak}}%
  {\newcommand{\pbreak}{\begin{center} \# \# \# \end{center}}}

% \renewcommand{\pbreak}{\thoughtbreak}
% Series Title Page
%==================
\def\seriestitle#1{\def\theseriestitle{\uppercase{#1}}}
\def\seriespage{\thispagestyle{empty}
\begin{center}
  \noindent{\large\theseriestitle}
  \vskip10.5pt
  \noindent Other Books in the series:
  \vskip4pt
  \theseriesbooks
\end{center}\cleardoublepage}
\def\seriesbook#1{\vskip.5pt{\noindent\hskip8pt\emph{#1}}\vskip1sp}
\def\seriesbooks#1\endseriesbooks{\def\theseriesbooks{\let\book\seriesbook{#1}}}

% Indices
%=========
\makeglossary
\makeindex
\let\oldindex\index{}%
\newcommand{\marked}[1]{\whenfinal{#1}{\underline{#1}}}%
\renewcommand{\index}[1]{\oldindex{#1}\marked{#1}}%
\newcommand{\Index}[2]{\oldindex{#1}\marked{#2}}%
\newcommand{\orbital}[1]{\oldindex{Orbital!#1}\marked{#1}}%
\newcommand{\equipment}[1]{\oldindex{Equipment!#1}\marked{#1}}%
\newcommand{\Equipment}[2]{\oldindex{Equipment!#1}\marked{#2}}%
\newcommand{\System}[2]{\oldindex{System!#1}\marked{#2}}%
\newcommand{\system}[1]{\oldindex{System!#1}\marked{#1}}%
\newcommand{\character}[1]{\oldindex{Character!#1}\marked{#1}}%
\newcommand{\Character}[2]{\oldindex{Character!#1}\marked{#2}}%
\newcommand{\ship}[1]{\oldindex{Ship!#1}\marked{\emph{#1}}}%
\newcommand{\Ship}[1]{\oldindex{Ship!#1}\marked{\emph{#1}}}%
\newcommand{\location}[1]{\oldindex{Location!#1}\marked{#1}}%
\newcommand{\Location}[2]{\oldindex{Location!#1}\marked{#2}}%
\newcommand{\clin}[1]{\oldindex{Clin!#1}Clin \marked{#1}}%

% Book-Specific Metadata
%========================

\def\NextBook{}
\def\ReturningCharacters{}
\title{carmilla}
\author{Benjamin Wilson}
\date{2016-10-10}
\newdate{firstprint}{10}{10}{2016}


\seriestitle{Series TItle}
\seriesbooks
  \book{Book in Series}
  \book{Book in Series}
\endseriesbooks


% Drafting Help
%==============
\newcommand{\tk}[1]{%
  \ifdefined\isdraft%
      % \textbf{TK-#1}
      \todo{#1}\textbf{TK-#1}%
    \else%
      \textbf{TK-#1}%
      % \todo{#1}\textbf{TK-#1}
  \fi%
}%
\newcommand{\TK}[1]{\tk{BAD-TK--#1}}

% Document
%========================
\begin{document}

% Frontmatter
%-------------
\frontmatter
  % r.1 - Half-Title
  \begin{titlingpage}
    \pagestyle{empty}
    \begin{center}
    \vspace*{2in}

    \Huge\textbf{\textsf\thetitle}

    \vspace*{0.75in}

    \Large\textsf\theauthor

    \vspace*{\fill}
    \end{center}
  \end{titlingpage}

  % r.2 - Series Title Page
  \seriespage
 \newpage
  % r.3 Titlepage
  \begin{titlingpage}
    \pagestyle{empty}
    \begin{center}
    \vspace*{\fill}

    \HUGE\textbf{\textsf\thetitle}

    \vspace*{0.25in}
    \line(1,0){150}
    \vspace*{0.25in}

    \Large\textsf\theauthor

    \vspace*{\fill}

    \vspace*{\fill}
    \includegraphics[width=2in]{_images/logo.png}\\[0cm]
    %\hspace*{\fill}\textsf{Dausha}\hspace*{\fill}\newline%
    %\textsf{Publishing}
    \end{center}
  \end{titlingpage}

  % r.4 Copyright Page
  \vspace*{\fill}
  \pagestyle{empty}

  \par\noindent\emph{\thetitle}
  \newline

  

  \par\noindent This is a work of fiction. Names, characters, places and incidents are either
  the product of the author's imagination or are used fictitiously, and any resemblance to
  actual persons, living or dead, business establishments, events or locales is entirely
  coincidental.\newline
  

  \par\noindent\emph{Copyright \copyright{} 2016 Benjamin Wilson.}\newline

  
  \par\noindent Cover Design by Benjamin Wilson\newline
  

  
  \par\noindent Book Design by Ben Wilson\newline
  

  


  \par\noindent All rights reserved.\newline

  \par\noindent No part of this publication may be reproduced, stored in a retrieval system, posted on the Internet, or transmitted, in any form or by any means, electronic, mechanical, photocopying, recording, or otherwise, without prior written permission from the author. The only exception is by a reviewer, who may quote short excerpts in a review.\newline

  \par\noindent \theauthor
  \par\noindent Visit my website at \url{http://example.com}\newline

  \par\noindent Printed in the United States of America
  \newline

  \par\noindent\textit{First Printing, \usmonthyear\displaydate{firstprint}}
  \newline

  \par\noindent ISBN-13 9-87654-3231-0123
  \vspace*{\fill}

  % r.5 - Dedication
  
  


  \newpage
\mainmatter
  \pagestyle{ebook}
  \sloppy

The Project Gutenberg EBook of Carmilla, by J. Sheridan LeFanu

This eBook is for the use of anyone anywhere at no cost and with
almost no restrictions whatsoever.  You may copy it, give it away or
re-use it under the terms of the Project Gutenberg License included
with this eBook or online at www.gutenberg.net

Title: Carmilla

Author: J. Sheridan LeFanu

Release Date: November 7, 2003 [EBook \#10007]
[Date last updated: December 1, 2004]

Language: English

\emph{** START OF THIS PROJECT GUTENBERG EBOOK CARMILLA **}

Produced by Suzanne Shell, Sjaani and PG Distributed Proofreaders

CARMILLA

J. Sheridan LeFanu

1872

\pbreak{}

PROLOGUE

\_Upon a paper attached to the Narrative which follows, Doctor Hesselius
has written a rather elaborate note, which he accompanies with a
reference to his Essay on the strange subject which the MS. illuminates.

This mysterious subject he treats, in that Essay, with his usual
learning and acumen, and with remarkable directness and condensation. It
will form but one volume of the series of that extraordinary man's
collected papers.

As I publish the case, in this volume, simply to interest the ``laity,'' I
shall forestall the intelligent lady, who relates it, in nothing; and
after due consideration, I have determined, therefore, to abstain from
presenting any précis of the learned Doctor's reasoning, or extract from
his statement on a subject which he describes as ``involving, not
improbably, some of the profoundest arcana of our dual existence, and
its intermediates.''

I was anxious on discovering this paper, to reopen the correspondence
commenced by Doctor Hesselius, so many years before, with a person so
clever and careful as his informant seems to have been. Much to my
regret, however, I found that she had died in the interval.

She, probably, could have added little to the Narrative \emph{which she
communicates in the following pages, with, so far as I can pronounce,
such conscientious particularity}.

\pbreak{}

\chapter{An Early Fright}\hypertarget{an-early-fright}{}\label{an-early-fright}

In Styria, we, though by no means magnificent people, inhabit a castle,
or schloss. A small income, in that part of the world, goes a great way.
Eight or nine hundred a year does wonders. Scantily enough ours would
have answered among wealthy people at home. My father is English, and I
bear an English name, although I never saw England. But here, in this
lonely and primitive place, where everything is so marvelously cheap, I
really don't see how ever so much more money would at all materially add
to our comforts, or even luxuries.

My father was in the Austrian service, and retired upon a pension and
his patrimony, and purchased this feudal residence, and the small estate
on which it stands, a bargain.

Nothing can be more picturesque or solitary. It stands on a slight
eminence in a forest. The road, very old and narrow, passes in front of
its drawbridge, never raised in my time, and its moat, stocked with
perch, and sailed over by many swans, and floating on its surface white
fleets of water lilies.

Over all this the schloss shows its many-windowed front; its towers,
and its Gothic chapel.

The forest opens in an irregular and very picturesque glade before its
gate, and at the right a steep Gothic bridge carries the road over a
stream that winds in deep shadow through the wood. I have said that this
is a very lonely place. Judge whether I say truth. Looking from the hall
door towards the road, the forest in which our castle stands extends
fifteen miles to the right, and twelve to the left. The nearest
inhabited village is about seven of your English miles to the left. The
nearest inhabited schloss of any historic associations, is that of old
General Spielsdorf, nearly twenty miles away to the right.

I have said ``the nearest \emph{inhabited} village,'' because there is, only
three miles westward, that is to say in the direction of General
Spielsdorf's schloss, a ruined village, with its quaint little church,
now roofless, in the aisle of which are the moldering tombs of the proud
family of Karnstein, now extinct, who once owned the equally desolate
chateau which, in the thick of the forest, overlooks the silent ruins
of the town.

Respecting the cause of the desertion of this striking and melancholy
spot, there is a legend which I shall relate to you another time.

I must tell you now, how very small is the party who constitute the
inhabitants of our castle. I don't include servants, or those dependents
who occupy rooms in the buildings attached to the schloss. Listen, and
wonder! My father, who is the kindest man on earth, but growing old; and
I, at the date of my story, only nineteen. Eight years have passed
since then.

I and my father constituted the family at the schloss. My mother, a
Styrian lady, died in my infancy, but I had a good-natured governess,
who had been with me from, I might almost say, my infancy. I could not
remember the time when her fat, benignant face was not a familiar
picture in my memory.

This was Madame Perrodon, a native of Berne, whose care and good nature
now in part supplied to me the loss of my mother, whom I do not even
remember, so early I lost her. She made a third at our little dinner
party. There was a fourth, Mademoiselle De Lafontaine, a lady such as
you term, I believe, a ``finishing governess.'' She spoke French and
German, Madame Perrodon French and broken English, to which my father
and I added English, which, partly to prevent its becoming a lost
language among us, and partly from patriotic motives, we spoke every
day. The consequence was a Babel, at which strangers used to laugh, and
which I shall make no attempt to reproduce in this narrative. And there
were two or three young lady friends besides, pretty nearly of my own
age, who were occasional visitors, for longer or shorter terms; and
these visits I sometimes returned.

These were our regular social resources; but of course there were chance
visits from ``neighbors'' of only five or six leagues distance. My life
was, notwithstanding, rather a solitary one, I can assure you.

My gouvernantes had just so much control over me as you might conjecture
such sage persons would have in the case of a rather spoiled girl, whose
only parent allowed her pretty nearly her own way in everything.

The first occurrence in my existence, which produced a terrible
impression upon my mind, which, in fact, never has been effaced, was one
of the very earliest incidents of my life which I can recollect. Some
people will think it so trifling that it should not be recorded here.
You will see, however, by-and-by, why I mention it. The nursery, as it
was called, though I had it all to myself, was a large room in the upper
story of the castle, with a steep oak roof. I can't have been more than
six years old, when one night I awoke, and looking round the room from
my bed, failed to see the nursery maid. Neither was my nurse there; and
I thought myself alone. I was not frightened, for I was one of those
happy children who are studiously kept in ignorance of ghost stories, of
fairy tales, and of all such lore as makes us cover up our heads when
the door cracks suddenly, or the flicker of an expiring candle makes the
shadow of a bedpost dance upon the wall, nearer to our faces. I was
vexed and insulted at finding myself, as I conceived, neglected, and I
began to whimper, preparatory to a hearty bout of roaring; when to my
surprise, I saw a solemn, but very pretty face looking at me from the
side of the bed. It was that of a young lady who was kneeling, with her
hands under the coverlet. I looked at her with a kind of pleased wonder,
and ceased whimpering. She caressed me with her hands, and lay down
beside me on the bed, and drew me towards her, smiling; I felt
immediately delightfully soothed, and fell asleep again. I was wakened
by a sensation as if two needles ran into my breast very deep at the
same moment, and I cried loudly. The lady started back, with her eyes
fixed on me, and then slipped down upon the floor, and, as I thought,
hid herself under the bed.

I was now for the first time frightened, and I yelled with all my might
and main. Nurse, nursery maid, housekeeper, all came running in, and
hearing my story, they made light of it, soothing me all they could
meanwhile. But, child as I was, I could perceive that their faces were
pale with an unwonted look of anxiety, and I saw them look under the
bed, and about the room, and peep under tables and pluck open cupboards;
and the housekeeper whispered to the nurse: ``Lay your hand along that
hollow in the bed; someone \emph{did} lie there, so sure as you did not; the
place is still warm.''

I remember the nursery maid petting me, and all three examining my
chest, where I told them I felt the puncture, and pronouncing that there
was no sign visible that any such thing had happened to me.

The housekeeper and the two other servants who were in charge of the
nursery, remained sitting up all night; and from that time a servant
always sat up in the nursery until I was about fourteen.

I was very nervous for a long time after this. A doctor was called in,
he was pallid and elderly. How well I remember his long saturnine face,
slightly pitted with smallpox, and his chestnut wig. For a good while,
every second day, he came and gave me medicine, which of course I hated.

The morning after I saw this apparition I was in a state of terror, and
could not bear to be left alone, daylight though it was, for a moment.

I remember my father coming up and standing at the bedside, and talking
cheerfully, and asking the nurse a number of questions, and laughing
very heartily at one of the answers; and patting me on the shoulder, and
kissing me, and telling me not to be frightened, that it was nothing but
a dream and could not hurt me.

But I was not comforted, for I knew the visit of the strange woman was
\emph{not} a dream; and I was \emph{awfully} frightened.

I was a little consoled by the nursery maid's assuring me that it was
she who had come and looked at me, and lain down beside me in the bed,
and that I must have been half-dreaming not to have known her face. But
this, though supported by the nurse, did not quite satisfy me.

I remembered, in the course of that day, a venerable old man, in a black
cassock, coming into the room with the nurse and housekeeper, and
talking a little to them, and very kindly to me; his face was very sweet
and gentle, and he told me they were going to pray, and joined my hands
together, and desired me to say, softly, while they were praying, ``Lord
hear all good prayers for us, for Jesus' sake.'' I think these were the
very words, for I often repeated them to myself, and my nurse used for
years to make me say them in my prayers.

I remembered so well the thoughtful sweet face of that white-haired old
man, in his black cassock, as he stood in that rude, lofty, brown room,
with the clumsy furniture of a fashion three hundred years old about
him, and the scanty light entering its shadowy atmosphere through the
small lattice. He kneeled, and the three women with him, and he prayed
aloud with an earnest quavering voice for, what appeared to me, a long
time. I forget all my life preceding that event, and for some time after
it is all obscure also, but the scenes I have just described stand out
vivid as the isolated pictures of the phantasmagoria surrounded
by darkness.

\pbreak{}

\chapter{A Guest}\hypertarget{a-guest}{}\label{a-guest}

I am now going to tell you something so strange that it will require all
your faith in my veracity to believe my story. It is not only true,
nevertheless, but truth of which I have been an eyewitness.

It was a sweet summer evening, and my father asked me, as he sometimes
did, to take a little ramble with him along that beautiful forest vista
which I have mentioned as lying in front of the schloss.

``General Spielsdorf cannot come to us so soon as I had hoped,'' said my
father, as we pursued our walk.

He was to have paid us a visit of some weeks, and we had expected his
arrival next day. He was to have brought with him a young lady, his
niece and ward, Mademoiselle Rheinfeldt, whom I had never seen, but whom
I had heard described as a very charming girl, and in whose society I
had promised myself many happy days. I was more disappointed than a
young lady living in a town, or a bustling neighborhood can possibly
imagine. This visit, and the new acquaintance it promised, had furnished
my day dream for many weeks.

``And how soon does he come?'' I asked.

``Not till autumn. Not for two months, I dare say,'' he answered. ``And I
am very glad now, dear, that you never knew Mademoiselle Rheinfeldt.''

``And why?'' I asked, both mortified and curious.

``Because the poor young lady is dead,'' he replied. ``I quite forgot I had
not told you, but you were not in the room when I received the General's
letter this evening.''

I was very much shocked. General Spielsdorf had mentioned in his first
letter, six or seven weeks before, that she was not so well as he would
wish her, but there was nothing to suggest the remotest suspicion
of danger.

``Here is the General's letter,'' he said, handing it to me. ``I am afraid
he is in great affliction; the letter appears to me to have been written
very nearly in distraction.''

We sat down on a rude bench, under a group of magnificent lime trees.
The sun was setting with all its melancholy splendor behind the sylvan
horizon, and the stream that flows beside our home, and passes under the
steep old bridge I have mentioned, wound through many a group of noble
trees, almost at our feet, reflecting in its current the fading crimson
of the sky. General Spielsdorf's letter was so extraordinary, so
vehement, and in some places so self-contradictory, that I read it twice
over--the second time aloud to my father--and was still unable to
account for it, except by supposing that grief had unsettled his mind.

It said ``I have lost my darling daughter, for as such I loved her.
During the last days of dear Bertha's illness I was not able to write
to you.

``Before then I had no idea of her danger. I have lost her, and now learn
\emph{all}, too late. She died in the peace of innocence, and in the glorious
hope of a blessed futurity. The fiend who betrayed our infatuated
hospitality has done it all. I thought I was receiving into my house
innocence, gaiety, a charming companion for my lost Bertha. Heavens!
what a fool have I been!

``I thank God my child died without a suspicion of the cause of her
sufferings. She is gone without so much as conjecturing the nature of
her illness, and the accursed passion of the agent of all this misery. I
devote my remaining days to tracking and extinguishing a monster. I am
told I may hope to accomplish my righteous and merciful purpose. At
present there is scarcely a gleam of light to guide me. I curse my
conceited incredulity, my despicable affectation of superiority, my
blindness, my obstinacy--all--too late. I cannot write or talk
collectedly now. I am distracted. So soon as I shall have a little
recovered, I mean to devote myself for a time to enquiry, which may
possibly lead me as far as Vienna. Some time in the autumn, two months
hence, or earlier if I live, I will see you--that is, if you permit me;
I will then tell you all that I scarce dare put upon paper now.
Farewell. Pray for me, dear friend.''

In these terms ended this strange letter. Though I had never seen Bertha
Rheinfeldt my eyes filled with tears at the sudden intelligence; I was
startled, as well as profoundly disappointed.

The sun had now set, and it was twilight by the time I had returned the
General's letter to my father.

It was a soft clear evening, and we loitered, speculating upon the
possible meanings of the violent and incoherent sentences which I had
just been reading. We had nearly a mile to walk before reaching the road
that passes the schloss in front, and by that time the moon was shining
brilliantly. At the drawbridge we met Madame Perrodon and Mademoiselle
De Lafontaine, who had come out, without their bonnets, to enjoy the
exquisite moonlight.

We heard their voices gabbling in animated dialogue as we approached. We
joined them at the drawbridge, and turned about to admire with them the
beautiful scene.

The glade through which we had just walked lay before us. At our left
the narrow road wound away under clumps of lordly trees, and was lost to
sight amid the thickening forest. At the right the same road crosses the
steep and picturesque bridge, near which stands a ruined tower which
once guarded that pass; and beyond the bridge an abrupt eminence rises,
covered with trees, and showing in the shadows some grey
ivy-clustered rocks.

Over the sward and low grounds a thin film of mist was stealing like
smoke, marking the distances with a transparent veil; and here and there
we could see the river faintly flashing in the moonlight.

No softer, sweeter scene could be imagined. The news I had just heard
made it melancholy; but nothing could disturb its character of profound
serenity, and the enchanted glory and vagueness of the prospect.

My father, who enjoyed the picturesque, and I, stood looking in silence
over the expanse beneath us. The two good governesses, standing a little
way behind us, discoursed upon the scene, and were eloquent upon
the moon.

Madame Perrodon was fat, middle-aged, and romantic, and talked and
sighed poetically. Mademoiselle De Lafontaine--in right of her father
who was a German, assumed to be psychological, metaphysical, and
something of a mystic--now declared that when the moon shone with a
light so intense it was well known that it indicated a special spiritual
activity. The effect of the full moon in such a state of brilliancy was
manifold. It acted on dreams, it acted on lunacy, it acted on nervous
people, it had marvelous physical influences connected with life.
Mademoiselle related that her cousin, who was mate of a merchant ship,
having taken a nap on deck on such a night, lying on his back, with his
face full in the light on the moon, had wakened, after a dream of an old
woman clawing him by the cheek, with his features horribly drawn to one
side; and his countenance had never quite recovered its equilibrium.

``The moon, this night,'' she said, ``is full of idyllic and magnetic
influence--and see, when you look behind you at the front of the schloss
how all its windows flash and twinkle with that silvery splendor, as if
unseen hands had lighted up the rooms to receive fairy guests.''

There are indolent styles of the spirits in which, indisposed to talk
ourselves, the talk of others is pleasant to our listless ears; and I
gazed on, pleased with the tinkle of the ladies' conversation.

``I have got into one of my moping moods tonight,'' said my father, after
a silence, and quoting Shakespeare, whom, by way of keeping up our
English, he used to read aloud, he said:

``{}`In truth I know not why I am so sad.
  It wearies me: you say it wearies you;
  But how I got it--came by it.'

``I forget the rest. But I feel as if some great misfortune were hanging
over us. I suppose the poor General's afflicted letter has had something
to do with it.''

At this moment the unwonted sound of carriage wheels and many hoofs upon
the road, arrested our attention.

They seemed to be approaching from the high ground overlooking the
bridge, and very soon the equipage emerged from that point. Two horsemen
first crossed the bridge, then came a carriage drawn by four horses, and
two men rode behind.

It seemed to be the traveling carriage of a person of rank; and we were
all immediately absorbed in watching that very unusual spectacle. It
became, in a few moments, greatly more interesting, for just as the
carriage had passed the summit of the steep bridge, one of the leaders,
taking fright, communicated his panic to the rest, and after a plunge or
two, the whole team broke into a wild gallop together, and dashing
between the horsemen who rode in front, came thundering along the road
towards us with the speed of a hurricane.

The excitement of the scene was made more painful by the clear,
long-drawn screams of a female voice from the carriage window.

We all advanced in curiosity and horror; me rather in silence, the rest
with various ejaculations of terror.

Our suspense did not last long. Just before you reach the castle
drawbridge, on the route they were coming, there stands by the roadside
a magnificent lime tree, on the other stands an ancient stone cross, at
sight of which the horses, now going at a pace that was perfectly
frightful, swerved so as to bring the wheel over the projecting roots
of the tree.

I knew what was coming. I covered my eyes, unable to see it out, and
turned my head away; at the same moment I heard a cry from my lady
friends, who had gone on a little.

Curiosity opened my eyes, and I saw a scene of utter confusion. Two of
the horses were on the ground, the carriage lay upon its side with two
wheels in the air; the men were busy removing the traces, and a lady
with a commanding air and figure had got out, and stood with clasped
hands, raising the handkerchief that was in them every now and then
to her eyes.

Through the carriage door was now lifted a young lady, who appeared to
be lifeless. My dear old father was already beside the elder lady, with
his hat in his hand, evidently tendering his aid and the resources of
his schloss. The lady did not appear to hear him, or to have eyes for
anything but the slender girl who was being placed against the slope
of the bank.

I approached; the young lady was apparently stunned, but she was
certainly not dead. My father, who piqued himself on being something of
a physician, had just had his fingers on her wrist and assured the lady,
who declared herself her mother, that her pulse, though faint and
irregular, was undoubtedly still distinguishable. The lady clasped her
hands and looked upward, as if in a momentary transport of gratitude;
but immediately she broke out again in that theatrical way which is, I
believe, natural to some people.

She was what is called a fine looking woman for her time of life, and
must have been handsome; she was tall, but not thin, and dressed in
black velvet, and looked rather pale, but with a proud and commanding
countenance, though now agitated strangely.

``Who was ever being so born to calamity?'' I heard her say, with clasped
hands, as I came up. ``Here am I, on a journey of life and death, in
prosecuting which to lose an hour is possibly to lose all. My child will
not have recovered sufficiently to resume her route for who can say how
long. I must leave her: I cannot, dare not, delay. How far on, sir, can
you tell, is the nearest village? I must leave her there; and shall not
see my darling, or even hear of her till my return, three months hence.''

I plucked my father by the coat, and whispered earnestly in his ear:
``Oh! papa, pray ask her to let her stay with us--it would be so
delightful. Do, pray.''

``If Madame will entrust her child to the care of my daughter, and of her
good gouvernante, Madame Perrodon, and permit her to remain as our
guest, under my charge, until her return, it will confer a distinction
and an obligation upon us, and we shall treat her with all the care and
devotion which so sacred a trust deserves.''

``I cannot do that, sir, it would be to task your kindness and chivalry
too cruelly,'' said the lady, distractedly.

``It would, on the contrary, be to confer on us a very great kindness at
the moment when we most need it. My daughter has just been disappointed
by a cruel misfortune, in a visit from which she had long anticipated a
great deal of happiness. If you confide this young lady to our care it
will be her best consolation. The nearest village on your route is
distant, and affords no such inn as you could think of placing your
daughter at; you cannot allow her to continue her journey for any
considerable distance without danger. If, as you say, you cannot suspend
your journey, you must part with her tonight, and nowhere could you do
so with more honest assurances of care and tenderness than here.''

There was something in this lady's air and appearance so distinguished
and even imposing, and in her manner so engaging, as to impress one,
quite apart from the dignity of her equipage, with a conviction that she
was a person of consequence.

By this time the carriage was replaced in its upright position, and the
horses, quite tractable, in the traces again.

The lady threw on her daughter a glance which I fancied was not quite so
affectionate as one might have anticipated from the beginning of the
scene; then she beckoned slightly to my father, and withdrew two or
three steps with him out of hearing; and talked to him with a fixed and
stern countenance, not at all like that with which she had
hitherto spoken.

I was filled with wonder that my father did not seem to perceive the
change, and also unspeakably curious to learn what it could be that she
was speaking, almost in his ear, with so much earnestness and rapidity.

Two or three minutes at most I think she remained thus employed, then
she turned, and a few steps brought her to where her daughter lay,
supported by Madame Perrodon. She kneeled beside her for a moment and
whispered, as Madame supposed, a little benediction in her ear; then
hastily kissing her she stepped into her carriage, the door was closed,
the footmen in stately liveries jumped up behind, the outriders spurred
on, the postilions cracked their whips, the horses plunged and broke
suddenly into a furious canter that threatened soon again to become a
gallop, and the carriage whirled away, followed at the same rapid pace
by the two horsemen in the rear.

\pbreak{}

\chapter{We Compare Notes}\hypertarget{we-compare-notes}{}\label{we-compare-notes}

We followed the \emph{cortege} with our eyes until it was swiftly lost to
sight in the misty wood; and the very sound of the hoofs and the wheels
died away in the silent night air.

Nothing remained to assure us that the adventure had not been an
illusion of a moment but the young lady, who just at that moment opened
her eyes. I could not see, for her face was turned from me, but she
raised her head, evidently looking about her, and I heard a very sweet
voice ask complainingly, ``Where is mamma?''

Our good Madame Perrodon answered tenderly, and added some comfortable
assurances.

I then heard her ask:

``Where am I? What is this place?'' and after that she said, ``I don't see
the carriage; and Matska, where is she?''

Madame answered all her questions in so far as she understood them; and
gradually the young lady remembered how the misadventure came about, and
was glad to hear that no one in, or in attendance on, the carriage was
hurt; and on learning that her mamma had left her here, till her return
in about three months, she wept.

I was going to add my consolations to those of Madame Perrodon when
Mademoiselle De Lafontaine placed her hand upon my arm, saying:

``Don't approach, one at a time is as much as she can at present converse
with; a very little excitement would possibly overpower her now.''

As soon as she is comfortably in bed, I thought, I will run up to her
room and see her.

My father in the meantime had sent a servant on horseback for the
physician, who lived about two leagues away; and a bedroom was being
prepared for the young lady's reception.

The stranger now rose, and leaning on Madame's arm, walked slowly over
the drawbridge and into the castle gate.

In the hall, servants waited to receive her, and she was conducted
forthwith to her room. The room we usually sat in as our drawing room is
long, having four windows, that looked over the moat and drawbridge,
upon the forest scene I have just described.

It is furnished in old carved oak, with large carved cabinets, and the
chairs are cushioned with crimson Utrecht velvet. The walls are covered
with tapestry, and surrounded with great gold frames, the figures being
as large as life, in ancient and very curious costume, and the subjects
represented are hunting, hawking, and generally festive. It is not too
stately to be extremely comfortable; and here we had our tea, for with
his usual patriotic leanings he insisted that the national beverage
should make its appearance regularly with our coffee and chocolate.

We sat here this night, and with candles lighted, were talking over the
adventure of the evening.

Madame Perrodon and Mademoiselle De Lafontaine were both of our party.
The young stranger had hardly lain down in her bed when she sank into a
deep sleep; and those ladies had left her in the care of a servant.

``How do you like our guest?'' I asked, as soon as Madame entered. ``Tell
me all about her?''

``I like her extremely,'' answered Madame, ``she is, I almost think, the
prettiest creature I ever saw; about your age, and so gentle and nice.''

``She is absolutely beautiful,'' threw in Mademoiselle, who had peeped for
a moment into the stranger's room.

``And such a sweet voice!'' added Madame Perrodon.

``Did you remark a woman in the carriage, after it was set up again, who
did not get out,'' inquired Mademoiselle, ``but only looked from
the window?''

``No, we had not seen her.''

Then she described a hideous black woman, with a sort of colored turban
on her head, and who was gazing all the time from the carriage window,
nodding and grinning derisively towards the ladies, with gleaming eyes
and large white eyeballs, and her teeth set as if in fury.

``Did you remark what an ill-looking pack of men the servants were?''
asked Madame.

``Yes,'' said my father, who had just come in, ``ugly, hang-dog looking
fellows as ever I beheld in my life. I hope they mayn't rob the poor
lady in the forest. They are clever rogues, however; they got everything
to rights in a minute.''

``I dare say they are worn out with too long traveling,'' said Madame.

``Besides looking wicked, their faces were so strangely lean, and dark,
and sullen. I am very curious, I own; but I dare say the young lady will
tell you all about it tomorrow, if she is sufficiently recovered.''

``I don't think she will,'' said my father, with a mysterious smile, and a
little nod of his head, as if he knew more about it than he cared
to tell us.

This made us all the more inquisitive as to what had passed between him
and the lady in the black velvet, in the brief but earnest interview
that had immediately preceded her departure.

We were scarcely alone, when I entreated him to tell me. He did not need
much pressing.

``There is no particular reason why I should not tell you. She expressed
a reluctance to trouble us with the care of her daughter, saying she was
in delicate health, and nervous, but not subject to any kind of
seizure--she volunteered that--nor to any illusion; being, in fact,
perfectly sane.''

``How very odd to say all that!'' I interpolated. ``It was so unnecessary.''

``At all events it \emph{was} said,'' he laughed, ``and as you wish to know all
that passed, which was indeed very little, I tell you. She then said, `I
am making a long journey of \emph{vital} importance--she emphasized the
word--rapid and secret; I shall return for my child in three months; in
the meantime, she will be silent as to who we are, whence we come, and
whither we are traveling.' That is all she said. She spoke very pure
French. When she said the word `secret,' she paused for a few seconds,
looking sternly, her eyes fixed on mine. I fancy she makes a great point
of that. You saw how quickly she was gone. I hope I have not done a very
foolish thing, in taking charge of the young lady.''

For my part, I was delighted. I was longing to see and talk to her; and
only waiting till the doctor should give me leave. You, who live in
towns, can have no idea how great an event the introduction of a new
friend is, in such a solitude as surrounded us.

The doctor did not arrive till nearly one o'clock; but I could no more
have gone to my bed and slept, than I could have overtaken, on foot, the
carriage in which the princess in black velvet had driven away.

When the physician came down to the drawing room, it was to report very
favorably upon his patient. She was now sitting up, her pulse quite
regular, apparently perfectly well. She had sustained no injury, and the
little shock to her nerves had passed away quite harmlessly. There could
be no harm certainly in my seeing her, if we both wished it; and, with
this permission I sent, forthwith, to know whether she would allow me to
visit her for a few minutes in her room.

The servant returned immediately to say that she desired nothing more.

You may be sure I was not long in availing myself of this permission.

Our visitor lay in one of the handsomest rooms in the schloss. It was,
perhaps, a little stately. There was a somber piece of tapestry opposite
the foot of the bed, representing Cleopatra with the asps to her bosom;
and other solemn classic scenes were displayed, a little faded, upon the
other walls. But there was gold carving, and rich and varied color
enough in the other decorations of the room, to more than redeem the
gloom of the old tapestry.

There were candles at the bedside. She was sitting up; her slender
pretty figure enveloped in the soft silk dressing gown, embroidered with
flowers, and lined with thick quilted silk, which her mother had thrown
over her feet as she lay upon the ground.

What was it that, as I reached the bedside and had just begun my little
greeting, struck me dumb in a moment, and made me recoil a step or two
from before her? I will tell you.

I saw the very face which had visited me in my childhood at night, which
remained so fixed in my memory, and on which I had for so many years so
often ruminated with horror, when no one suspected of what I
was thinking.

It was pretty, even beautiful; and when I first beheld it, wore the
same melancholy expression.

But this almost instantly lighted into a strange fixed smile of
recognition.

There was a silence of fully a minute, and then at length she spoke; I
could not.

``How wonderful!'' she exclaimed. ``Twelve years ago, I saw your face in a
dream, and it has haunted me ever since.''

``Wonderful indeed!'' I repeated, overcoming with an effort the horror
that had for a time suspended my utterances. ``Twelve years ago, in
vision or reality, I certainly saw you. I could not forget your face. It
has remained before my eyes ever since.''

Her smile had softened. Whatever I had fancied strange in it, was gone,
and it and her dimpling cheeks were now delightfully pretty and
intelligent.

I felt reassured, and continued more in the vein which hospitality
indicated, to bid her welcome, and to tell her how much pleasure her
accidental arrival had given us all, and especially what a happiness it
was to me.

I took her hand as I spoke. I was a little shy, as lonely people are,
but the situation made me eloquent, and even bold. She pressed my hand,
she laid hers upon it, and her eyes glowed, as, looking hastily into
mine, she smiled again, and blushed.

She answered my welcome very prettily. I sat down beside her, still
wondering; and she said:

``I must tell you my vision about you; it is so very strange that you and
I should have had, each of the other so vivid a dream, that each should
have seen, I you and you me, looking as we do now, when of course we
both were mere children. I was a child, about six years old, and I awoke
from a confused and troubled dream, and found myself in a room, unlike
my nursery, wainscoted clumsily in some dark wood, and with cupboards
and bedsteads, and chairs, and benches placed about it. The beds were,
I thought, all empty, and the room itself without anyone but myself in
it; and I, after looking about me for some time, and admiring especially
an iron candlestick with two branches, which I should certainly know
again, crept under one of the beds to reach the window; but as I got
from under the bed, I heard someone crying; and looking up, while I was
still upon my knees, I saw you--most assuredly you--as I see you now; a
beautiful young lady, with golden hair and large blue eyes, and
lips--your lips--you as you are here.

``Your looks won me; I climbed on the bed and put my arms about you, and
I think we both fell asleep. I was aroused by a scream; you were sitting
up screaming. I was frightened, and slipped down upon the ground, and,
it seemed to me, lost consciousness for a moment; and when I came to
myself, I was again in my nursery at home. Your face I have never
forgotten since. I could not be misled by mere resemblance. \emph{You are}
the lady whom I saw then.''

It was now my turn to relate my corresponding vision, which I did, to
the undisguised wonder of my new acquaintance.

``I don't know which should be most afraid of the other,'' she said, again
smiling--``If you were less pretty I think I should be very much afraid
of you, but being as you are, and you and I both so young, I feel only
that I have made your acquaintance twelve years ago, and have already a
right to your intimacy; at all events it does seem as if we were
destined, from our earliest childhood, to be friends. I wonder whether
you feel as strangely drawn towards me as I do to you; I have never had
a friend--shall I find one now?'' She sighed, and her fine dark eyes
gazed passionately on me.

Now the truth is, I felt rather unaccountably towards the beautiful
stranger. I did feel, as she said, ``drawn towards her,'' but there was
also something of repulsion. In this ambiguous feeling, however, the
sense of attraction immensely prevailed. She interested and won me; she
was so beautiful and so indescribably engaging.

I perceived now something of languor and exhaustion stealing over her,
and hastened to bid her good night.

``The doctor thinks,'' I added, ``that you ought to have a maid to sit up
with you tonight; one of ours is waiting, and you will find her a very
useful and quiet creature.''

``How kind of you, but I could not sleep, I never could with an attendant
in the room. I shan't require any assistance--and, shall I confess my
weakness, I am haunted with a terror of robbers. Our house was robbed
once, and two servants murdered, so I always lock my door. It has become
a habit--and you look so kind I know you will forgive me. I see there is
a key in the lock.''

She held me close in her pretty arms for a moment and whispered in my
ear, ``Good night, darling, it is very hard to part with you, but good
night; tomorrow, but not early, I shall see you again.''

She sank back on the pillow with a sigh, and her fine eyes followed me
with a fond and melancholy gaze, and she murmured again ``Good night,
dear friend.''

Young people like, and even love, on impulse. I was flattered by the
evident, though as yet undeserved, fondness she showed me. I liked the
confidence with which she at once received me. She was determined that
we should be very near friends.

Next day came and we met again. I was delighted with my companion; that
is to say, in many respects.

Her looks lost nothing in daylight--she was certainly the most beautiful
creature I had ever seen, and the unpleasant remembrance of the face
presented in my early dream, had lost the effect of the first unexpected
recognition.

She confessed that she had experienced a similar shock on seeing me, and
precisely the same faint antipathy that had mingled with my admiration
of her. We now laughed together over our momentary horrors.

\pbreak{}

\chapter{Her Habits--A Saunter}\hypertarget{her-habits--a-saunter}{}\label{her-habits--a-saunter}

I told you that I was charmed with her in most particulars.

There were some that did not please me so well.

She was above the middle height of women. I shall begin by describing
her.

She was slender, and wonderfully graceful. Except that her movements
were languid--very languid--indeed, there was nothing in her appearance
to indicate an invalid. Her complexion was rich and brilliant; her
features were small and beautifully formed; her eyes large, dark, and
lustrous; her hair was quite wonderful, I never saw hair so
magnificently thick and long when it was down about her shoulders; I
have often placed my hands under it, and laughed with wonder at its
weight. It was exquisitely fine and soft, and in color a rich very dark
brown, with something of gold. I loved to let it down, tumbling with its
own weight, as, in her room, she lay back in her chair talking in her
sweet low voice, I used to fold and braid it, and spread it out and
play with it. Heavens! If I had but known all!

I said there were particulars which did not please me. I have told you
that her confidence won me the first night I saw her; but I found that
she exercised with respect to herself, her mother, her history,
everything in fact connected with her life, plans, and people, an ever
wakeful reserve. I dare say I was unreasonable, perhaps I was wrong; I
dare say I ought to have respected the solemn injunction laid upon my
father by the stately lady in black velvet. But curiosity is a restless
and unscrupulous passion, and no one girl can endure, with patience,
that hers should be baffled by another. What harm could it do anyone to
tell me what I so ardently desired to know? Had she no trust in my good
sense or honor? Why would she not believe me when I assured her, so
solemnly, that I would not divulge one syllable of what she told me to
any mortal breathing.

There was a coldness, it seemed to me, beyond her years, in her smiling
melancholy persistent refusal to afford me the least ray of light.

I cannot say we quarreled upon this point, for she would not quarrel
upon any. It was, of course, very unfair of me to press her, very
ill-bred, but I really could not help it; and I might just as well have
let it alone.

What she did tell me amounted, in my unconscionable estimation--to
nothing.

It was all summed up in three very vague disclosures:

First--Her name was Carmilla.

Second--Her family was very ancient and noble.

Third--Her home lay in the direction of the west.

She would not tell me the name of her family, nor their armorial
bearings, nor the name of their estate, nor even that of the country
they lived in.

You are not to suppose that I worried her incessantly on these subjects.
I watched opportunity, and rather insinuated than urged my inquiries.
Once or twice, indeed, I did attack her more directly. But no matter
what my tactics, utter failure was invariably the result. Reproaches and
caresses were all lost upon her. But I must add this, that her evasion
was conducted with so pretty a melancholy and deprecation, with so many,
and even passionate declarations of her liking for me, and trust in my
honor, and with so many promises that I should at last know all, that I
could not find it in my heart long to be offended with her.

She used to place her pretty arms about my neck, draw me to her, and
laying her cheek to mine, murmur with her lips near my ear, ``Dearest,
your little heart is wounded; think me not cruel because I obey the
irresistible law of my strength and weakness; if your dear heart is
wounded, my wild heart bleeds with yours. In the rapture of my enormous
humiliation I live in your warm life, and you shall die--die, sweetly
die--into mine. I cannot help it; as I draw near to you, you, in your
turn, will draw near to others, and learn the rapture of that cruelty,
which yet is love; so, for a while, seek to know no more of me and mine,
but trust me with all your loving spirit.''

And when she had spoken such a rhapsody, she would press me more closely
in her trembling embrace, and her lips in soft kisses gently glow
upon my cheek.

Her agitations and her language were unintelligible to me.

From these foolish embraces, which were not of very frequent occurrence,
I must allow, I used to wish to extricate myself; but my energies seemed
to fail me. Her murmured words sounded like a lullaby in my ear, and
soothed my resistance into a trance, from which I only seemed to recover
myself when she withdrew her arms.

In these mysterious moods I did not like her. I experienced a strange
tumultuous excitement that was pleasurable, ever and anon, mingled with
a vague sense of fear and disgust. I had no distinct thoughts about her
while such scenes lasted, but I was conscious of a love growing into
adoration, and also of abhorrence. This I know is paradox, but I can
make no other attempt to explain the feeling.

I now write, after an interval of more than ten years, with a trembling
hand, with a confused and horrible recollection of certain occurrences
and situations, in the ordeal through which I was unconsciously passing;
though with a vivid and very sharp remembrance of the main current of
my story.

But, I suspect, in all lives there are certain emotional scenes, those
in which our passions have been most wildly and terribly roused, that
are of all others the most vaguely and dimly remembered.

Sometimes after an hour of apathy, my strange and beautiful companion
would take my hand and hold it with a fond pressure, renewed again and
again; blushing softly, gazing in my face with languid and burning eyes,
and breathing so fast that her dress rose and fell with the tumultuous
respiration. It was like the ardor of a lover; it embarrassed me; it was
hateful and yet over-powering; and with gloating eyes she drew me to
her, and her hot lips traveled along my cheek in kisses; and she would
whisper, almost in sobs, ``You are mine, you \emph{shall} be mine, you and I
are one for ever.'' Then she had thrown herself back in her chair, with
her small hands over her eyes, leaving me trembling.

``Are we related,'' I used to ask; ``what can you mean by all this? I
remind you perhaps of someone whom you love; but you must not, I hate
it; I don't know you--I don't know myself when you look so and talk so.''

She used to sigh at my vehemence, then turn away and drop my hand.

Respecting these very extraordinary manifestations I strove in vain to
form any satisfactory theory--I could not refer them to affectation or
trick. It was unmistakably the momentary breaking out of suppressed
instinct and emotion. Was she, notwithstanding her mother's volunteered
denial, subject to brief visitations of insanity; or was there here a
disguise and a romance? I had read in old storybooks of such things.
What if a boyish lover had found his way into the house, and sought to
prosecute his suit in masquerade, with the assistance of a clever old
adventuress. But there were many things against this hypothesis, highly
interesting as it was to my vanity.

I could boast of no little attentions such as masculine gallantry
delights to offer. Between these passionate moments there were long
intervals of commonplace, of gaiety, of brooding melancholy, during
which, except that I detected her eyes so full of melancholy fire,
following me, at times I might have been as nothing to her. Except in
these brief periods of mysterious excitement her ways were girlish; and
there was always a languor about her, quite incompatible with a
masculine system in a state of health.

In some respects her habits were odd. Perhaps not so singular in the
opinion of a town lady like you, as they appeared to us rustic people.
She used to come down very late, generally not till one o'clock, she
would then take a cup of chocolate, but eat nothing; we then went out
for a walk, which was a mere saunter, and she seemed, almost
immediately, exhausted, and either returned to the schloss or sat on one
of the benches that were placed, here and there, among the trees. This
was a bodily languor in which her mind did not sympathize. She was
always an animated talker, and very intelligent.

She sometimes alluded for a moment to her own home, or mentioned an
adventure or situation, or an early recollection, which indicated a
people of strange manners, and described customs of which we knew
nothing. I gathered from these chance hints that her native country was
much more remote than I had at first fancied.

As we sat thus one afternoon under the trees a funeral passed us by. It
was that of a pretty young girl, whom I had often seen, the daughter of
one of the rangers of the forest. The poor man was walking behind the
coffin of his darling; she was his only child, and he looked quite
heartbroken.

Peasants walking two-and-two came behind, they were singing a funeral
hymn.

I rose to mark my respect as they passed, and joined in the hymn they
were very sweetly singing.

My companion shook me a little roughly, and I turned surprised.

She said brusquely, ``Don't you perceive how discordant that is?''

``I think it very sweet, on the contrary,'' I answered, vexed at the
interruption, and very uncomfortable, lest the people who composed the
little procession should observe and resent what was passing.

I resumed, therefore, instantly, and was again interrupted. ``You pierce
my ears,'' said Carmilla, almost angrily, and stopping her ears with her
tiny fingers. ``Besides, how can you tell that your religion and mine are
the same; your forms wound me, and I hate funerals. What a fuss! Why you
must die--\emph{everyone} must die; and all are happier when they do.
Come home.''

``My father has gone on with the clergyman to the churchyard. I thought
you knew she was to be buried today.''

``She? I don't trouble my head about peasants. I don't know who she is,''
answered Carmilla, with a flash from her fine eyes.

``She is the poor girl who fancied she saw a ghost a fortnight ago, and
has been dying ever since, till yesterday, when she expired.''

``Tell me nothing about ghosts. I shan't sleep tonight if you do.''

``I hope there is no plague or fever coming; all this looks very like
it,'' I continued. ``The swineherd's young wife died only a week ago, and
she thought something seized her by the throat as she lay in her bed,
and nearly strangled her. Papa says such horrible fancies do accompany
some forms of fever. She was quite well the day before. She sank
afterwards, and died before a week.''

``Well, \emph{her} funeral is over, I hope, and \emph{her} hymn sung; and our ears
shan't be tortured with that discord and jargon. It has made me nervous.
Sit down here, beside me; sit close; hold my hand; press it
hard-hard-harder.''

We had moved a little back, and had come to another seat.

She sat down. Her face underwent a change that alarmed and even
terrified me for a moment. It darkened, and became horribly livid; her
teeth and hands were clenched, and she frowned and compressed her lips,
while she stared down upon the ground at her feet, and trembled all over
with a continued shudder as irrepressible as ague. All her energies
seemed strained to suppress a fit, with which she was then breathlessly
tugging; and at length a low convulsive cry of suffering broke from her,
and gradually the hysteria subsided. ``There! That comes of strangling
people with hymns!'' she said at last. ``Hold me, hold me still. It is
passing away.''

And so gradually it did; and perhaps to dissipate the somber impression
which the spectacle had left upon me, she became unusually animated and
chatty; and so we got home.

This was the first time I had seen her exhibit any definable symptoms of
that delicacy of health which her mother had spoken of. It was the first
time, also, I had seen her exhibit anything like temper.

Both passed away like a summer cloud; and never but once afterwards did
I witness on her part a momentary sign of anger. I will tell you how
it happened.

She and I were looking out of one of the long drawing room windows, when
there entered the courtyard, over the drawbridge, a figure of a wanderer
whom I knew very well. He used to visit the schloss generally twice
a year.

It was the figure of a hunchback, with the sharp lean features that
generally accompany deformity. He wore a pointed black beard, and he was
smiling from ear to ear, showing his white fangs. He was dressed in
buff, black, and scarlet, and crossed with more straps and belts than I
could count, from which hung all manner of things. Behind, he carried a
magic lantern, and two boxes, which I well knew, in one of which was a
salamander, and in the other a mandrake. These monsters used to make my
father laugh. They were compounded of parts of monkeys, parrots,
squirrels, fish, and hedgehogs, dried and stitched together with great
neatness and startling effect. He had a fiddle, a box of conjuring
apparatus, a pair of foils and masks attached to his belt, several other
mysterious cases dangling about him, and a black staff with copper
ferrules in his hand. His companion was a rough spare dog, that followed
at his heels, but stopped short, suspiciously at the drawbridge, and in
a little while began to howl dismally.

In the meantime, the mountebank, standing in the midst of the courtyard,
raised his grotesque hat, and made us a very ceremonious bow, paying his
compliments very volubly in execrable French, and German not
much better.

Then, disengaging his fiddle, he began to scrape a lively air to which
he sang with a merry discord, dancing with ludicrous airs and activity,
that made me laugh, in spite of the dog's howling.

Then he advanced to the window with many smiles and salutations, and
his hat in his left hand, his fiddle under his arm, and with a fluency
that never took breath, he gabbled a long advertisement of all his
accomplishments, and the resources of the various arts which he placed
at our service, and the curiosities and entertainments which it was in
his power, at our bidding, to display.

``Will your ladyships be pleased to buy an amulet against the oupire,
which is going like the wolf, I hear, through these woods,'' he said
dropping his hat on the pavement. ``They are dying of it right and left
and here is a charm that never fails; only pinned to the pillow, and you
may laugh in his face.''

These charms consisted of oblong slips of vellum, with cabalistic
ciphers and diagrams upon them.

Carmilla instantly purchased one, and so did I.

He was looking up, and we were smiling down upon him, amused; at least,
I can answer for myself. His piercing black eye, as he looked up in our
faces, seemed to detect something that fixed for a moment his curiosity.
In an instant he unrolled a leather case, full of all manner of odd
little steel instruments.

``See here, my lady,'' he said, displaying it, and addressing me, ``I
profess, among other things less useful, the art of dentistry. Plague
take the dog!'' he interpolated. ``Silence, beast! He howls so that your
ladyships can scarcely hear a word. Your noble friend, the young lady at
your right, has the sharpest tooth,--long, thin, pointed, like an awl,
like a needle; ha, ha! With my sharp and long sight, as I look up, I
have seen it distinctly; now if it happens to hurt the young lady, and I
think it must, here am I, here are my file, my punch, my nippers; I will
make it round and blunt, if her ladyship pleases; no longer the tooth of
a fish, but of a beautiful young lady as she is. Hey? Is the young lady
displeased? Have I been too bold? Have I offended her?''

The young lady, indeed, looked very angry as she drew back from the
window.

``How dares that mountebank insult us so? Where is your father? I shall
demand redress from him. My father would have had the wretch tied up to
the pump, and flogged with a cart whip, and burnt to the bones with the
cattle brand!''

She retired from the window a step or two, and sat down, and had hardly
lost sight of the offender, when her wrath subsided as suddenly as it
had risen, and she gradually recovered her usual tone, and seemed to
forget the little hunchback and his follies.

My father was out of spirits that evening. On coming in he told us that
there had been another case very similar to the two fatal ones which had
lately occurred. The sister of a young peasant on his estate, only a
mile away, was very ill, had been, as she described it, attacked very
nearly in the same way, and was now slowly but steadily sinking.

``All this,'' said my father, ``is strictly referable to natural causes.
These poor people infect one another with their superstitions, and so
repeat in imagination the images of terror that have infested their
neighbors.''

``But that very circumstance frightens one horribly,'' said Carmilla.

``How so?'' inquired my father.

``I am so afraid of fancying I see such things; I think it would be as
bad as reality.''

``We are in God's hands: nothing can happen without his permission, and
all will end well for those who love him. He is our faithful creator; He
has made us all, and will take care of us.''

``Creator! \emph{Nature!}'' said the young lady in answer to my gentle father.
``And this disease that invades the country is natural. Nature. All
things proceed from Nature--don't they? All things in the heaven, in the
earth, and under the earth, act and live as Nature ordains? I
think so.''

``The doctor said he would come here today,'' said my father, after a
silence. ``I want to know what he thinks about it, and what he thinks we
had better do.''

``Doctors never did me any good,'' said Carmilla.

``Then you have been ill?'' I asked.

``More ill than ever you were,'' she answered.

``Long ago?''

``Yes, a long time. I suffered from this very illness; but I forget all
but my pain and weakness, and they were not so bad as are suffered in
other diseases.''

``You were very young then?''

``I dare say, let us talk no more of it. You would not wound a friend?''

She looked languidly in my eyes, and passed her arm round my waist
lovingly, and led me out of the room. My father was busy over some
papers near the window.

``Why does your papa like to frighten us?'' said the pretty girl with a
sigh and a little shudder.

``He doesn't, dear Carmilla, it is the very furthest thing from his
mind.''

``Are you afraid, dearest?''

``I should be very much if I fancied there was any real danger of my
being attacked as those poor people were.''

``You are afraid to die?''

``Yes, every one is.''

``But to die as lovers may--to die together, so that they may live
together.

``Girls are caterpillars while they live in the world, to be finally
butterflies when the summer comes; but in the meantime there are grubs
and larvae, don't you see--each with their peculiar propensities,
necessities and structure. So says Monsieur Buffon, in his big book, in
the next room.''

Later in the day the doctor came, and was closeted with papa for some
time.

He was a skilful man, of sixty and upwards, he wore powder, and shaved
his pale face as smooth as a pumpkin. He and papa emerged from the room
together, and I heard papa laugh, and say as they came out:

``Well, I do wonder at a wise man like you. What do you say to
hippogriffs and dragons?''

The doctor was smiling, and made answer, shaking his head--

``Nevertheless life and death are mysterious states, and we know little
of the resources of either.''

And so they walked on, and I heard no more. I did not then know what the
doctor had been broaching, but I think I guess it now.

\pbreak{}

\chapter{A Wonderful Likeness}\hypertarget{a-wonderful-likeness}{}\label{a-wonderful-likeness}

This evening there arrived from Gratz the grave, dark-faced son of the
picture cleaner, with a horse and cart laden with two large packing
cases, having many pictures in each. It was a journey of ten leagues,
and whenever a messenger arrived at the schloss from our little capital
of Gratz, we used to crowd about him in the hall, to hear the news.

This arrival created in our secluded quarters quite a sensation. The
cases remained in the hall, and the messenger was taken charge of by the
servants till he had eaten his supper. Then with assistants, and armed
with hammer, ripping chisel, and turnscrew, he met us in the hall, where
we had assembled to witness the unpacking of the cases.

Carmilla sat looking listlessly on, while one after the other the old
pictures, nearly all portraits, which had undergone the process of
renovation, were brought to light. My mother was of an old Hungarian
family, and most of these pictures, which were about to be restored to
their places, had come to us through her.

My father had a list in his hand, from which he read, as the artist
rummaged out the corresponding numbers. I don't know that the pictures
were very good, but they were, undoubtedly, very old, and some of them
very curious also. They had, for the most part, the merit of being now
seen by me, I may say, for the first time; for the smoke and dust of
time had all but obliterated them.

``There is a picture that I have not seen yet,'' said my father. ``In one
corner, at the top of it, is the name, as well as I could read, `Marcia
Karnstein,' and the date `1698'; and I am curious to see how it has
turned out.''

I remembered it; it was a small picture, about a foot and a half high,
and nearly square, without a frame; but it was so blackened by age that
I could not make it out.

The artist now produced it, with evident pride. It was quite beautiful;
it was startling; it seemed to live. It was the effigy of Carmilla!

``Carmilla, dear, here is an absolute miracle. Here you are, living,
smiling, ready to speak, in this picture. Isn't it beautiful, Papa? And
see, even the little mole on her throat.''

My father laughed, and said ``Certainly it is a wonderful likeness,'' but
he looked away, and to my surprise seemed but little struck by it, and
went on talking to the picture cleaner, who was also something of an
artist, and discoursed with intelligence about the portraits or other
works, which his art had just brought into light and color, while I was
more and more lost in wonder the more I looked at the picture.

``Will you let me hang this picture in my room, papa?'' I asked.

``Certainly, dear,'' said he, smiling, ``I'm very glad you think it so
like. It must be prettier even than I thought it, if it is.''

The young lady did not acknowledge this pretty speech, did not seem to
hear it. She was leaning back in her seat, her fine eyes under their
long lashes gazing on me in contemplation, and she smiled in a kind
of rapture.

``And now you can read quite plainly the name that is written in the
corner. It is not Marcia; it looks as if it was done in gold. The name
is Mircalla, Countess Karnstein, and this is a little coronet over and
underneath A.D. 1698. I am descended from the Karnsteins; that is,
mamma was.''

``Ah!'' said the lady, languidly, ``so am I, I think, a very long descent,
very ancient. Are there any Karnsteins living now?''

``None who bear the name, I believe. The family were ruined, I believe,
in some civil wars, long ago, but the ruins of the castle are only about
three miles away.''

``How interesting!'' she said, languidly. ``But see what beautiful
moonlight!'' She glanced through the hall door, which stood a little
open. ``Suppose you take a little ramble round the court, and look down
at the road and river.''

``It is so like the night you came to us,'' I said.

She sighed; smiling.

She rose, and each with her arm about the other's waist, we walked out
upon the pavement.

In silence, slowly we walked down to the drawbridge, where the beautiful
landscape opened before us.

``And so you were thinking of the night I came here?'' she almost
whispered.

``Are you glad I came?''

``Delighted, dear Carmilla,'' I answered.

``And you asked for the picture you think like me, to hang in your room,''
she murmured with a sigh, as she drew her arm closer about my waist, and
let her pretty head sink upon my shoulder. ``How romantic you are,
Carmilla,'' I said. ``Whenever you tell me your story, it will be made up
chiefly of some one great romance.''

She kissed me silently.

``I am sure, Carmilla, you have been in love; that there is, at this
moment, an affair of the heart going on.''

``I have been in love with no one, and never shall,'' she whispered,
``unless it should be with you.''

How beautiful she looked in the moonlight!

Shy and strange was the look with which she quickly hid her face in my
neck and hair, with tumultuous sighs, that seemed almost to sob, and
pressed in mine a hand that trembled.

Her soft cheek was glowing against mine. ``Darling, darling,'' she
murmured, ``I live in you; and you would die for me, I love you so.''

I started from her.

She was gazing on me with eyes from which all fire, all meaning had
flown, and a face colorless and apathetic.

``Is there a chill in the air, dear?'' she said drowsily. ``I almost
shiver; have I been dreaming? Let us come in. Come; come; come in.''

``You look ill, Carmilla; a little faint. You certainly must take some
wine,'' I said.

``Yes. I will. I'm better now. I shall be quite well in a few minutes.
Yes, do give me a little wine,'' answered Carmilla, as we approached
the door.

``Let us look again for a moment; it is the last time, perhaps, I shall
see the moonlight with you.''

``How do you feel now, dear Carmilla? Are you really better?'' I asked.

I was beginning to take alarm, lest she should have been stricken with
the strange epidemic that they said had invaded the country about us.

``Papa would be grieved beyond measure,'' I added, ``if he thought you were
ever so little ill, without immediately letting us know. We have a very
skilful doctor near us, the physician who was with papa today.''

``I'm sure he is. I know how kind you all are; but, dear child, I am
quite well again. There is nothing ever wrong with me, but a
little weakness.

``People say I am languid; I am incapable of exertion; I can scarcely walk
as far as a child of three years old: and every now and then the little
strength I have falters, and I become as you have just seen me. But
after all I am very easily set up again; in a moment I am perfectly
myself. See how I have recovered.''

So, indeed, she had; and she and I talked a great deal, and very
animated she was; and the remainder of that evening passed without any
recurrence of what I called her infatuations. I mean her crazy talk and
looks, which embarrassed, and even frightened me.

But there occurred that night an event which gave my thoughts quite a
new turn, and seemed to startle even Carmilla's languid nature into
momentary energy.

\pbreak{}

\chapter{A Very Strange Agony}\hypertarget{a-very-strange-agony}{}\label{a-very-strange-agony}

When we got into the drawing room, and had sat down to our coffee and
chocolate, although Carmilla did not take any, she seemed quite herself
again, and Madame, and Mademoiselle De Lafontaine, joined us, and made a
little card party, in the course of which papa came in for what he
called his ``dish of tea.''

When the game was over he sat down beside Carmilla on the sofa, and
asked her, a little anxiously, whether she had heard from her mother
since her arrival.

She answered ``No.''

He then asked whether she knew where a letter would reach her at
present.

``I cannot tell,'' she answered ambiguously, ``but I have been thinking of
leaving you; you have been already too hospitable and too kind to me. I
have given you an infinity of trouble, and I should wish to take a
carriage tomorrow, and post in pursuit of her; I know where I shall
ultimately find her, although I dare not yet tell you.''

``But you must not dream of any such thing,'' exclaimed my father, to my
great relief. ``We can't afford to lose you so, and I won't consent to
your leaving us, except under the care of your mother, who was so good
as to consent to your remaining with us till she should herself return.
I should be quite happy if I knew that you heard from her: but this
evening the accounts of the progress of the mysterious disease that has
invaded our neighborhood, grow even more alarming; and my beautiful
guest, I do feel the responsibility, unaided by advice from your mother,
very much. But I shall do my best; and one thing is certain, that you
must not think of leaving us without her distinct direction to that
effect. We should suffer too much in parting from you to consent to
it easily.''

``Thank you, sir, a thousand times for your hospitality,'' she answered,
smiling bashfully. ``You have all been too kind to me; I have seldom been
so happy in all my life before, as in your beautiful chateau, under your
care, and in the society of your dear daughter.''

So he gallantly, in his old-fashioned way, kissed her hand, smiling and
pleased at her little speech.

I accompanied Carmilla as usual to her room, and sat and chatted with
her while she was preparing for bed.

``Do you think,'' I said at length, ``that you will ever confide fully in
me?''

She turned round smiling, but made no answer, only continued to smile on
me.

``You won't answer that?'' I said. ``You can't answer pleasantly; I ought
not to have asked you.''

``You were quite right to ask me that, or anything. You do not know how
dear you are to me, or you could not think any confidence too great to
look for.  But I am under vows, no nun half so awfully, and I dare not
tell my story yet, even to you. The time is very near when you shall
know everything. You will think me cruel, very selfish, but love is
always selfish; the more ardent the more selfish. How jealous I am you
cannot know. You must come with me, loving me, to death; or else hate me
and still come with me, and \emph{hating} me through death and after. There
is no such word as indifference in my apathetic nature.''

``Now, Carmilla, you are going to talk your wild nonsense again,'' I said
hastily.

``Not I, silly little fool as I am, and full of whims and fancies; for
your sake I'll talk like a sage. Were you ever at a ball?''

``No; how you do run on. What is it like? How charming it must be.''

``I almost forget, it is years ago.''

I laughed.

``You are not so old. Your first ball can hardly be forgotten yet.''

``I remember everything about it--with an effort. I see it all, as divers
see what is going on above them, through a medium, dense, rippling, but
transparent. There occurred that night what has confused the picture,
and made its colours faint. I was all but assassinated in my bed,
wounded here,'' she touched her breast, ``and never was the same since.''

``Were you near dying?''

``Yes, very--a cruel love--strange love, that would have taken my life.
Love will have its sacrifices. No sacrifice without blood. Let us go to
sleep now; I feel so lazy. How can I get up just now and lock my door?''

She was lying with her tiny hands buried in her rich wavy hair, under
her cheek, her little head upon the pillow, and her glittering eyes
followed me wherever I moved, with a kind of shy smile that I could
not decipher.

I bid her good night, and crept from the room with an uncomfortable
sensation.

I often wondered whether our pretty guest ever said her prayers. I
certainly had never seen her upon her knees. In the morning she never
came down until long after our family prayers were over, and at night
she never left the drawing room to attend our brief evening prayers
in the hall.

If it had not been that it had casually come out in one of our careless
talks that she had been baptised, I should have doubted her being a
Christian. Religion was a subject on which I had never heard her speak a
word. If I had known the world better, this particular neglect or
antipathy would not have so much surprised me.

The precautions of nervous people are infectious, and persons of a like
temperament are pretty sure, after a time, to imitate them. I had
adopted Carmilla's habit of locking her bedroom door, having taken into
my head all her whimsical alarms about midnight invaders and prowling
assassins. I had also adopted her precaution of making a brief search
through her room, to satisfy herself that no lurking assassin or robber
was ``ensconced.''

These wise measures taken, I got into my bed and fell asleep. A light
was burning in my room. This was an old habit, of very early date, and
which nothing could have tempted me to dispense with.

Thus fortified I might take my rest in peace. But dreams come through
stone walls, light up dark rooms, or darken light ones, and their
persons make their exits and their entrances as they please, and laugh
at locksmiths.

I had a dream that night that was the beginning of a very strange agony.

I cannot call it a nightmare, for I was quite conscious of being asleep.

But I was equally conscious of being in my room, and lying in bed,
precisely as I actually was. I saw, or fancied I saw, the room and its
furniture just as I had seen it last, except that it was very dark, and
I saw something moving round the foot of the bed, which at first I
could not accurately distinguish. But I soon saw that it was a
sooty-black animal that resembled a monstrous cat. It appeared to me
about four or five feet long for it measured fully the length of the
hearthrug as it passed over it; and it continued to-ing and fro-ing with
the lithe, sinister restlessness of a beast in a cage. I could not cry
out, although as you may suppose, I was terrified. Its pace was growing
faster, and the room rapidly darker and darker, and at length so dark
that I could no longer see anything of it but its eyes. I felt it spring
lightly on the bed. The two broad eyes approached my face, and suddenly
I felt a stinging pain as if two large needles darted, an inch or two
apart, deep into my breast. I waked with a scream. The room was lighted
by the candle that burnt there all through the night, and I saw a female
figure standing at the foot of the bed, a little at the right side. It
was in a dark loose dress, and its hair was down and covered its
shoulders. A block of stone could not have been more still. There was
not the slightest stir of respiration. As I stared at it, the figure
appeared to have changed its place, and was now nearer the door; then,
close to it, the door opened, and it passed out.

I was now relieved, and able to breathe and move. My first thought was
that Carmilla had been playing me a trick, and that I had forgotten to
secure my door. I hastened to it, and found it locked as usual on the
inside. I was afraid to open it--I was horrified. I sprang into my bed
and covered my head up in the bedclothes, and lay there more dead than
alive till morning.

\pbreak{}

\chapter{Descending}\hypertarget{descending}{}\label{descending}

It would be vain my attempting to tell you the horror with which, even
now, I recall the occurrence of that night. It was no such transitory
terror as a dream leaves behind it. It seemed to deepen by time, and
communicated itself to the room and the very furniture that had
encompassed the apparition.

I could not bear next day to be alone for a moment. I should have told
papa, but for two opposite reasons. At one time I thought he would laugh
at my story, and I could not bear its being treated as a jest; and at
another I thought he might fancy that I had been attacked by the
mysterious complaint which had invaded our neighborhood. I had myself no
misgiving of the kind, and as he had been rather an invalid for some
time, I was afraid of alarming him.

I was comfortable enough with my good-natured companions, Madame
Perrodon, and the vivacious Mademoiselle Lafontaine. They both perceived
that I was out of spirits and nervous, and at length I told them what
lay so heavy at my heart.

Mademoiselle laughed, but I fancied that Madame Perrodon looked anxious.

``By-the-by,'' said Mademoiselle, laughing, ``the long lime tree walk,
behind Carmilla's bedroom window, is haunted!''

``Nonsense!'' exclaimed Madame, who probably thought the theme rather
inopportune, ``and who tells that story, my dear?''

``Martin says that he came up twice, when the old yard gate was being
repaired, before sunrise, and twice saw the same female figure walking
down the lime tree avenue.''

``So he well might, as long as there are cows to milk in the river
fields,'' said Madame.

``I daresay; but Martin chooses to be frightened, and never did I see
fool more frightened.''

``You must not say a word about it to Carmilla, because she can see down
that walk from her room window,'' I interposed, ``and she is, if possible,
a greater coward than I.''

Carmilla came down rather later than usual that day.

``I was so frightened last night,'' she said, so soon as were together,
``and I am sure I should have seen something dreadful if it had not been
for that charm I bought from the poor little hunchback whom I called
such hard names. I had a dream of something black coming round my bed,
and I awoke in a perfect horror, and I really thought, for some seconds,
I saw a dark figure near the chimney-piece, but I felt under my pillow
for my charm, and the moment my fingers touched it, the figure
disappeared, and I felt quite certain, only that I had it by me, that
something frightful would have made its appearance, and, perhaps,
throttled me, as it did those poor people we heard of.

``Well, listen to me,'' I began, and recounted my adventure, at the
recital of which she appeared horrified.

``And had you the charm near you?'' she asked, earnestly.

``No, I had dropped it into a china vase in the drawing room, but I shall
certainly take it with me tonight, as you have so much faith in it.''

At this distance of time I cannot tell you, or even understand, how I
overcame my horror so effectually as to lie alone in my room that night.
I remember distinctly that I pinned the charm to my pillow. I fell
asleep almost immediately, and slept even more soundly than usual
all night.

Next night I passed as well. My sleep was delightfully deep and
dreamless.

But I wakened with a sense of lassitude and melancholy, which, however,
did not exceed a degree that was almost luxurious.

``Well, I told you so,'' said Carmilla, when I described my quiet sleep,
``I had such delightful sleep myself last night; I pinned the charm to
the breast of my nightdress. It was too far away the night before. I am
quite sure it was all fancy, except the dreams. I used to think that
evil spirits made dreams, but our doctor told me it is no such thing.
Only a fever passing by, or some other malady, as they often do, he
said, knocks at the door, and not being able to get in, passes on, with
that alarm.''

``And what do you think the charm is?'' said I.

``It has been fumigated or immersed in some drug, and is an antidote
against the malaria,'' she answered.

``Then it acts only on the body?''

``Certainly; you don't suppose that evil spirits are frightened by bits
of ribbon, or the perfumes of a druggist's shop? No, these complaints,
wandering in the air, begin by trying the nerves, and so infect the
brain, but before they can seize upon you, the antidote repels them.
That I am sure is what the charm has done for us. It is nothing magical,
it is simply natural.''

I should have been happier if I could have quite agreed with Carmilla,
but I did my best, and the impression was a little losing its force.

For some nights I slept profoundly; but still every morning I felt the
same lassitude, and a languor weighed upon me all day. I felt myself a
changed girl. A strange melancholy was stealing over me, a melancholy
that I would not have interrupted. Dim thoughts of death began to open,
and an idea that I was slowly sinking took gentle, and, somehow, not
unwelcome, possession of me. If it was sad, the tone of mind which this
induced was also sweet.

Whatever it might be, my soul acquiesced in it.

I would not admit that I was ill, I would not consent to tell my papa,
or to have the doctor sent for.

Carmilla became more devoted to me than ever, and her strange paroxysms
of languid adoration more frequent. She used to gloat on me with
increasing ardor the more my strength and spirits waned. This always
shocked me like a momentary glare of insanity.

Without knowing it, I was now in a pretty advanced stage of the
strangest illness under which mortal ever suffered. There was an
unaccountable fascination in its earlier symptoms that more than
reconciled me to the incapacitating effect of that stage of the malady.
This fascination increased for a time, until it reached a certain point,
when gradually a sense of the horrible mingled itself with it,
deepening, as you shall hear, until it discolored and perverted the
whole state of my life.

The first change I experienced was rather agreeable. It was very near
the turning point from which began the descent of Avernus.

Certain vague and strange sensations visited me in my sleep. The
prevailing one was of that pleasant, peculiar cold thrill which we feel
in bathing, when we move against the current of a river. This was soon
accompanied by dreams that seemed interminable, and were so vague that
I could never recollect their scenery and persons, or any one connected
portion of their action. But they left an awful impression, and a sense
of exhaustion, as if I had passed through a long period of great mental
exertion and danger.

After all these dreams there remained on waking a remembrance of having
been in a place very nearly dark, and of having spoken to people whom I
could not see; and especially of one clear voice, of a female's, very
deep, that spoke as if at a distance, slowly, and producing always the
same sensation of indescribable solemnity and fear. Sometimes there came
a sensation as if a hand was drawn softly along my cheek and neck.
Sometimes it was as if warm lips kissed me, and longer and longer and
more lovingly as they reached my throat, but there the caress fixed
itself. My heart beat faster, my breathing rose and fell rapidly and
full drawn; a sobbing, that rose into a sense of strangulation,
supervened, and turned into a dreadful convulsion, in which my senses
left me and I became unconscious.

It was now three weeks since the commencement of this unaccountable
state.

My sufferings had, during the last week, told upon my appearance. I had
grown pale, my eyes were dilated and darkened underneath, and the
languor which I had long felt began to display itself in my countenance.

My father asked me often whether I was ill; but, with an obstinacy which
now seems to me unaccountable, I persisted in assuring him that I was
quite well.

In a sense this was true. I had no pain, I could complain of no bodily
derangement. My complaint seemed to be one of the imagination, or the
nerves, and, horrible as my sufferings were, I kept them, with a morbid
reserve, very nearly to myself.

It could not be that terrible complaint which the peasants called the
oupire, for I had now been suffering for three weeks, and they were
seldom ill for much more than three days, when death put an end to
their miseries.

Carmilla complained of dreams and feverish sensations, but by no means
of so alarming a kind as mine. I say that mine were extremely alarming.
Had I been capable of comprehending my condition, I would have invoked
aid and advice on my knees. The narcotic of an unsuspected influence was
acting upon me, and my perceptions were benumbed.

I am going to tell you now of a dream that led immediately to an odd
discovery.

One night, instead of the voice I was accustomed to hear in the dark, I
heard one, sweet and tender, and at the same time terrible, which said,
``Your mother warns you to beware of the assassin.'' At the same time a
light unexpectedly sprang up, and I saw Carmilla, standing, near the
foot of my bed, in her white nightdress, bathed, from her chin to her
feet, in one great stain of blood.

I wakened with a shriek, possessed with the one idea that Carmilla was
being murdered. I remember springing from my bed, and my next
recollection is that of standing on the lobby, crying for help.

Madame and Mademoiselle came scurrying out of their rooms in alarm; a
lamp burned always on the lobby, and seeing me, they soon learned the
cause of my terror.

I insisted on our knocking at Carmilla's door. Our knocking was
unanswered.

It soon became a pounding and an uproar. We shrieked her name, but all
was vain.

We all grew frightened, for the door was locked. We hurried back, in
panic, to my room. There we rang the bell long and furiously. If my
father's room had been at that side of the house, we would have called
him up at once to our aid. But, alas! he was quite out of hearing, and
to reach him involved an excursion for which we none of us had courage.

Servants, however, soon came running up the stairs; I had got on my
dressing gown and slippers meanwhile, and my companions were already
similarly furnished. Recognizing the voices of the servants on the
lobby, we sallied out together; and having renewed, as fruitlessly, our
summons at Carmilla's door, I ordered the men to force the lock. They
did so, and we stood, holding our lights aloft, in the doorway, and so
stared into the room.

We called her by name; but there was still no reply. We looked round the
room. Everything was undisturbed. It was exactly in the state in which I
had left it on bidding her good night. But Carmilla was gone.

\pbreak{}

\chapter{Search}\hypertarget{search}{}\label{search}

At sight of the room, perfectly undisturbed except for our violent
entrance, we began to cool a little, and soon recovered our senses
sufficiently to dismiss the men. It had struck Mademoiselle that
possibly Carmilla had been wakened by the uproar at her door, and in her
first panic had jumped from her bed, and hid herself in a press, or
behind a curtain, from which she could not, of course, emerge until the
majordomo and his myrmidons had withdrawn. We now recommenced our
search, and began to call her name again.

It was all to no purpose. Our perplexity and agitation increased. We
examined the windows, but they were secured. I implored of Carmilla, if
she had concealed herself, to play this cruel trick no longer--to come
out and to end our anxieties. It was all useless. I was by this time
convinced that she was not in the room, nor in the dressing room, the
door of which was still locked on this side. She could not have passed
it. I was utterly puzzled. Had Carmilla discovered one of those secret
passages which the old housekeeper said were known to exist in the
schloss, although the tradition of their exact situation had been lost?
A little time would, no doubt, explain all--utterly perplexed as, for
the present, we were.

It was past four o'clock, and I preferred passing the remaining hours of
darkness in Madame's room. Daylight brought no solution of the
difficulty.

The whole household, with my father at its head, was in a state of
agitation next morning. Every part of the chateau was searched. The
grounds were explored. No trace of the missing lady could be discovered.
The stream was about to be dragged; my father was in distraction; what a
tale to have to tell the poor girl's mother on her return. I, too, was
almost beside myself, though my grief was quite of a different kind.

The morning was passed in alarm and excitement. It was now one o'clock,
and still no tidings. I ran up to Carmilla's room, and found her
standing at her dressing table. I was astounded. I could not believe my
eyes. She beckoned me to her with her pretty finger, in silence. Her
face expressed extreme fear.

I ran to her in an ecstasy of joy; I kissed and embraced her again and
again. I ran to the bell and rang it vehemently, to bring others to the
spot who might at once relieve my father's anxiety.

``Dear Carmilla, what has become of you all this time? We have been in
agonies of anxiety about you,'' I exclaimed. ``Where have you been? How
did you come back?''

``Last night has been a night of wonders,'' she said.

``For mercy's sake, explain all you can.''

``It was past two last night,'' she said, ``when I went to sleep as usual
in my bed, with my doors locked, that of the dressing room, and that
opening upon the gallery. My sleep was uninterrupted, and, so far as I
know, dreamless; but I woke just now on the sofa in the dressing room
there, and I found the door between the rooms open, and the other door
forced. How could all this have happened without my being wakened? It
must have been accompanied with a great deal of noise, and I am
particularly easily wakened; and how could I have been carried out of my
bed without my sleep having been interrupted, I whom the slightest stir
startles?''

By this time, Madame, Mademoiselle, my father, and a number of the
servants were in the room. Carmilla was, of course, overwhelmed with
inquiries, congratulations, and welcomes. She had but one story to tell,
and seemed the least able of all the party to suggest any way of
accounting for what had happened.

My father took a turn up and down the room, thinking. I saw Carmilla's
eye follow him for a moment with a sly, dark glance.

When my father had sent the servants away, Mademoiselle having gone in
search of a little bottle of valerian and salvolatile, and there being
no one now in the room with Carmilla, except my father, Madame, and
myself, he came to her thoughtfully, took her hand very kindly, led her
to the sofa, and sat down beside her.

``Will you forgive me, my dear, if I risk a conjecture, and ask a
question?''

``Who can have a better right?'' she said. ``Ask what you please, and I
will tell you everything. But my story is simply one of bewilderment and
darkness. I know absolutely nothing. Put any question you please, but
you know, of course, the limitations mamma has placed me under.''

``Perfectly, my dear child. I need not approach the topics on which she
desires our silence. Now, the marvel of last night consists in your
having been removed from your bed and your room, without being wakened,
and this removal having occurred apparently while the windows were still
secured, and the two doors locked upon the inside. I will tell you my
theory and ask you a question.''

Carmilla was leaning on her hand dejectedly; Madame and I were
listening breathlessly.

``Now, my question is this. Have you ever been suspected of walking in
your sleep?''

``Never, since I was very young indeed.''

``But you did walk in your sleep when you were young?''

``Yes; I know I did. I have been told so often by my old nurse.''

My father smiled and nodded.

``Well, what has happened is this. You got up in your sleep, unlocked the
door, not leaving the key, as usual, in the lock, but taking it out and
locking it on the outside; you again took the key out, and carried it
away with you to some one of the five-and-twenty rooms on this floor, or
perhaps upstairs or downstairs. There are so many rooms and closets, so
much heavy furniture, and such accumulations of lumber, that it would
require a week to search this old house thoroughly. Do you see, now,
what I mean?''

``I do, but not all,'' she answered.

``And how, papa, do you account for her finding herself on the sofa in
the dressing room, which we had searched so carefully?''

``She came there after you had searched it, still in her sleep, and at
last awoke spontaneously, and was as much surprised to find herself
where she was as any one else. I wish all mysteries were as easily and
innocently explained as yours, Carmilla,'' he said, laughing. ``And so we
may congratulate ourselves on the certainty that the most natural
explanation of the occurrence is one that involves no drugging, no
tampering with locks, no burglars, or poisoners, or witches--nothing
that need alarm Carmilla, or anyone else, for our safety.''

Carmilla was looking charmingly. Nothing could be more beautiful than
her tints. Her beauty was, I think, enhanced by that graceful languor
that was peculiar to her. I think my father was silently contrasting her
looks with mine, for he said:

``I wish my poor Laura was looking more like herself''; and he sighed.

So our alarms were happily ended, and Carmilla restored to her friends.

\pbreak{}

\chapter{The Doctor}\hypertarget{the-doctor}{}\label{the-doctor}

As Carmilla would not hear of an attendant sleeping in her room, my
father arranged that a servant should sleep outside her door, so that
she would not attempt to make another such excursion without being
arrested at her own door.

That night passed quietly; and next morning early, the doctor, whom my
father had sent for without telling me a word about it, arrived to
see me.

Madame accompanied me to the library; and there the grave little doctor,
with white hair and spectacles, whom I mentioned before, was waiting to
receive me.

I told him my story, and as I proceeded he grew graver and graver.

We were standing, he and I, in the recess of one of the windows, facing
one another. When my statement was over, he leaned with his shoulders
against the wall, and with his eyes fixed on me earnestly, with an
interest in which was a dash of horror.

After a minute's reflection, he asked Madame if he could see my father.

He was sent for accordingly, and as he entered, smiling, he said:

``I dare say, doctor, you are going to tell me that I am an old fool for
having brought you here; I hope I am.''

But his smile faded into shadow as the doctor, with a very grave face,
beckoned him to him.

He and the doctor talked for some time in the same recess where I had
just conferred with the physician. It seemed an earnest and
argumentative conversation. The room is very large, and I and Madame
stood together, burning with curiosity, at the farther end. Not a word
could we hear, however, for they spoke in a very low tone, and the deep
recess of the window quite concealed the doctor from view, and very
nearly my father, whose foot, arm, and shoulder only could we see; and
the voices were, I suppose, all the less audible for the sort of closet
which the thick wall and window formed.

After a time my father's face looked into the room; it was pale,
thoughtful, and, I fancied, agitated.

``Laura, dear, come here for a moment. Madame, we shan't trouble you, the
doctor says, at present.''

Accordingly I approached, for the first time a little alarmed; for,
although I felt very weak, I did not feel ill; and strength, one always
fancies, is a thing that may be picked up when we please.

My father held out his hand to me, as I drew near, but he was looking at
the doctor, and he said:

``It certainly is very odd; I don't understand it quite. Laura, come
here, dear; now attend to Doctor Spielsberg, and recollect yourself.''

``You mentioned a sensation like that of two needles piercing the skin,
somewhere about your neck, on the night when you experienced your first
horrible dream. Is there still any soreness?''

``None at all,'' I answered.

``Can you indicate with your finger about the point at which you think
this occurred?''

``Very little below my throat--here,'' I answered.

I wore a morning dress, which covered the place I pointed to.

``Now you can satisfy yourself,'' said the doctor. ``You won't mind your
papa's lowering your dress a very little. It is necessary, to detect a
symptom of the complaint under which you have been suffering.''

I acquiesced. It was only an inch or two below the edge of my collar.

``God bless me!--so it is,'' exclaimed my father, growing pale.

``You see it now with your own eyes,'' said the doctor, with a gloomy
triumph.

``What is it?'' I exclaimed, beginning to be frightened.

``Nothing, my dear young lady, but a small blue spot, about the size of
the tip of your little finger; and now,'' he continued, turning to papa,
``the question is what is best to be done?''

``Is there any danger?'' I urged, in great trepidation.

``I trust not, my dear,'' answered the doctor. ``I don't see why you should
not recover. I don't see why you should not begin immediately to get
better. That is the point at which the sense of strangulation begins?''

``Yes,'' I answered.

``And--recollect as well as you can--the same point was a kind of center
of that thrill which you described just now, like the current of a cold
stream running against you?''

``It may have been; I think it was.''

``Ay, you see?'' he added, turning to my father. ``Shall I say a word to
Madame?''

``Certainly,'' said my father.

He called Madame to him, and said:

``I find my young friend here far from well. It won't be of any great
consequence, I hope; but it will be necessary that some steps be taken,
which I will explain by-and-by; but in the meantime, Madame, you will
be so good as not to let Miss Laura be alone for one moment. That is the
only direction I need give for the present. It is indispensable.''

``We may rely upon your kindness, Madame, I know,'' added my father.

Madame satisfied him eagerly.

``And you, dear Laura, I know you will observe the doctor's direction.''

``I shall have to ask your opinion upon another patient, whose symptoms
slightly resemble those of my daughter, that have just been detailed to
you--very much milder in degree, but I believe quite of the same sort.
She is a young lady--our guest; but as you say you will be passing this
way again this evening, you can't do better than take your supper here,
and you can then see her. She does not come down till the afternoon.''

``I thank you,'' said the doctor. ``I shall be with you, then, at about
seven this evening.''

And then they repeated their directions to me and to Madame, and with
this parting charge my father left us, and walked out with the doctor;
and I saw them pacing together up and down between the road and the
moat, on the grassy platform in front of the castle, evidently absorbed
in earnest conversation.

The doctor did not return. I saw him mount his horse there, take his
leave, and ride away eastward through the forest.

Nearly at the same time I saw the man arrive from Dranfield with the
letters, and dismount and hand the bag to my father.

In the meantime, Madame and I were both busy, lost in conjecture as to
the reasons of the singular and earnest direction which the doctor and
my father had concurred in imposing. Madame, as she afterwards told me,
was afraid the doctor apprehended a sudden seizure, and that, without
prompt assistance, I might either lose my life in a fit, or at least be
seriously hurt.

The interpretation did not strike me; and I fancied, perhaps luckily for
my nerves, that the arrangement was prescribed simply to secure a
companion, who would prevent my taking too much exercise, or eating
unripe fruit, or doing any of the fifty foolish things to which young
people are supposed to be prone.

About half an hour after my father came in--he had a letter in his
hand--and said:

``This letter had been delayed; it is from General Spielsdorf. He might
have been here yesterday, he may not come till tomorrow or he may be
here today.''

He put the open letter into my hand; but he did not look pleased, as he
used when a guest, especially one so much loved as the General,
was coming.

On the contrary, he looked as if he wished him at the bottom of the Red
Sea. There was plainly something on his mind which he did not choose
to divulge.

``Papa, darling, will you tell me this?'' said I, suddenly laying my hand
on his arm, and looking, I am sure, imploringly in his face.

``Perhaps,'' he answered, smoothing my hair caressingly over my eyes.

``Does the doctor think me very ill?''

``No, dear; he thinks, if right steps are taken, you will be quite well
again, at least, on the high road to a complete recovery, in a day or
two,'' he answered, a little dryly. ``I wish our good friend, the General,
had chosen any other time; that is, I wish you had been perfectly well
to receive him.''

``But do tell me, papa,'' I insisted, ``what does he think is the matter
with me?''

``Nothing; you must not plague me with questions,'' he answered, with more
irritation than I ever remember him to have displayed before; and seeing
that I looked wounded, I suppose, he kissed me, and added, ``You shall
know all about it in a day or two; that is, all that I know. In the
meantime you are not to trouble your head about it.''

He turned and left the room, but came back before I had done wondering
and puzzling over the oddity of all this; it was merely to say that he
was going to Karnstein, and had ordered the carriage to be ready at
twelve, and that I and Madame should accompany him; he was going to see the
priest who lived near those picturesque grounds, upon business, and as
Carmilla had never seen them, she could follow, when she came down, with
Mademoiselle, who would bring materials for what you call a picnic,
which might be laid for us in the ruined castle.

At twelve o'clock, accordingly, I was ready, and not long after, my
father, Madame and I set out upon our projected drive.

Passing the drawbridge we turn to the right, and follow the road over
the steep Gothic bridge, westward, to reach the deserted village and
ruined castle of Karnstein.

No sylvan drive can be fancied prettier. The ground breaks into gentle
hills and hollows, all clothed with beautiful wood, totally destitute of
the comparative formality which artificial planting and early culture
and pruning impart.

The irregularities of the ground often lead the road out of its course,
and cause it to wind beautifully round the sides of broken hollows and
the steeper sides of the hills, among varieties of ground almost
inexhaustible.

Turning one of these points, we suddenly encountered our old friend, the
General, riding towards us, attended by a mounted servant. His
portmanteaus were following in a hired wagon, such as we term a cart.

The General dismounted as we pulled up, and, after the usual greetings,
was easily persuaded to accept the vacant seat in the carriage and send
his horse on with his servant to the schloss.

\pbreak{}

\chapter{Bereaved}\hypertarget{bereaved}{}\label{bereaved}

It was about ten months since we had last seen him: but that time had
sufficed to make an alteration of years in his appearance. He had grown
thinner; something of gloom and anxiety had taken the place of that
cordial serenity which used to characterize his features. His dark blue
eyes, always penetrating, now gleamed with a sterner light from under
his shaggy grey eyebrows. It was not such a change as grief alone
usually induces, and angrier passions seemed to have had their share in
bringing it about.

We had not long resumed our drive, when the General began to talk, with
his usual soldierly directness, of the bereavement, as he termed it,
which he had sustained in the death of his beloved niece and ward; and
he then broke out in a tone of intense bitterness and fury, inveighing
against the ``hellish arts'' to which she had fallen a victim, and
expressing, with more exasperation than piety, his wonder that Heaven
should tolerate so monstrous an indulgence of the lusts and malignity
of hell.

My father, who saw at once that something very extraordinary had
befallen, asked him, if not too painful to him, to detail the
circumstances which he thought justified the strong terms in which he
expressed himself.

``I should tell you all with pleasure,'' said the General, ``but you would
not believe me.''

``Why should I not?'' he asked.

``Because,'' he answered testily, ``you believe in nothing but what
consists with your own prejudices and illusions. I remember when I was
like you, but I have learned better.''

``Try me,'' said my father; ``I am not such a dogmatist as you suppose.
Besides which, I very well know that you generally require proof for
what you believe, and am, therefore, very strongly predisposed to
respect your conclusions.''

``You are right in supposing that I have not been led lightly into a
belief in the marvelous--for what I have experienced is marvelous--and I
have been forced by extraordinary evidence to credit that which ran
counter, diametrically, to all my theories. I have been made the dupe of
a preternatural conspiracy.''

Notwithstanding his professions of confidence in the General's
penetration, I saw my father, at this point, glance at the General,
with, as I thought, a marked suspicion of his sanity.

The General did not see it, luckily. He was looking gloomily and
curiously into the glades and vistas of the woods that were opening
before us.

``You are going to the Ruins of Karnstein?'' he said. ``Yes, it is a lucky
coincidence; do you know I was going to ask you to bring me there to
inspect them. I have a special object in exploring. There is a ruined
chapel, ain't there, with a great many tombs of that extinct family?''

``So there are--highly interesting,'' said my father. ``I hope you are
thinking of claiming the title and estates?''

My father said this gaily, but the General did not recollect the laugh,
or even the smile, which courtesy exacts for a friend's joke; on the
contrary, he looked grave and even fierce, ruminating on a matter that
stirred his anger and horror.

``Something very different,'' he said, gruffly. ``I mean to unearth some of
those fine people. I hope, by God's blessing, to accomplish a pious
sacrilege here, which will relieve our earth of certain monsters, and
enable honest people to sleep in their beds without being assailed by
murderers. I have strange things to tell you, my dear friend, such as I
myself would have scouted as incredible a few months since.''

My father looked at him again, but this time not with a glance of
suspicion--with an eye, rather, of keen intelligence and alarm.

``The house of Karnstein,'' he said, ``has been long extinct: a hundred
years at least. My dear wife was maternally descended from the
Karnsteins. But the name and title have long ceased to exist. The castle
is a ruin; the very village is deserted; it is fifty years since the
smoke of a chimney was seen there; not a roof left.''

``Quite true. I have heard a great deal about that since I last saw you;
a great deal that will astonish you. But I had better relate everything
in the order in which it occurred,'' said the General. ``You saw my dear
ward--my child, I may call her. No creature could have been more
beautiful, and only three months ago none more blooming.''

``Yes, poor thing! when I saw her last she certainly was quite lovely,''
said my father. ``I was grieved and shocked more than I can tell you, my
dear friend; I knew what a blow it was to you.''

He took the General's hand, and they exchanged a kind pressure. Tears
gathered in the old soldier's eyes. He did not seek to conceal them.
He said:

``We have been very old friends; I knew you would feel for me, childless
as I am. She had become an object of very near interest to me, and
repaid my care by an affection that cheered my home and made my life
happy. That is all gone. The years that remain to me on earth may not be
very long; but by God's mercy I hope to accomplish a service to mankind
before I die, and to subserve the vengeance of Heaven upon the fiends
who have murdered my poor child in the spring of her hopes and beauty!''

``You said, just now, that you intended relating everything as it
occurred,'' said my father. ``Pray do; I assure you that it is not mere
curiosity that prompts me.''

By this time we had reached the point at which the Drunstall road, by
which the General had come, diverges from the road which we were
traveling to Karnstein.

``How far is it to the ruins?'' inquired the General, looking anxiously
forward.

``About half a league,'' answered my father. ``Pray let us hear the story
you were so good as to promise.''

\pbreak{}

\chapter{The Story}\hypertarget{the-story}{}\label{the-story}

``With all my heart,'' said the General, with an effort; and after a short
pause in which to arrange his subject, he commenced one of the strangest
narratives I ever heard.

``My dear child was looking forward with great pleasure to the visit you
had been so good as to arrange for her to your charming daughter.'' Here
he made me a gallant but melancholy bow. ``In the meantime we had an
invitation to my old friend the Count Carlsfeld, whose schloss is about
six leagues to the other side of Karnstein. It was to attend the series
of fetes which, you remember, were given by him in honor of his
illustrious visitor, the Grand Duke Charles.''

``Yes; and very splendid, I believe, they were,'' said my father.

``Princely! But then his hospitalities are quite regal. He has Aladdin's
lamp. The night from which my sorrow dates was devoted to a magnificent
masquerade. The grounds were thrown open, the trees hung with colored
lamps. There was such a display of fireworks as Paris itself had never
witnessed. And such music--music, you know, is my weakness--such
ravishing music! The finest instrumental band, perhaps, in the world,
and the finest singers who could be collected from all the great operas
in Europe. As you wandered through these fantastically illuminated
grounds, the moon-lighted chateau throwing a rosy light from its long
rows of windows, you would suddenly hear these ravishing voices stealing
from the silence of some grove, or rising from boats upon the lake. I
felt myself, as I looked and listened, carried back into the romance and
poetry of my early youth.

``When the fireworks were ended, and the ball beginning, we returned to
the noble suite of rooms that were thrown open to the dancers. A masked
ball, you know, is a beautiful sight; but so brilliant a spectacle of
the kind I never saw before.

``It was a very aristocratic assembly. I was myself almost the only
`nobody' present.

``My dear child was looking quite beautiful. She wore no mask. Her
excitement and delight added an unspeakable charm to her features,
always lovely. I remarked a young lady, dressed magnificently, but
wearing a mask, who appeared to me to be observing my ward with
extraordinary interest. I had seen her, earlier in the evening, in the
great hall, and again, for a few minutes, walking near us, on the
terrace under the castle windows, similarly employed. A lady, also
masked, richly and gravely dressed, and with a stately air, like a
person of rank, accompanied her as a chaperon.

``Had the young lady not worn a mask, I could, of course, have been much
more certain upon the question whether she was really watching my
poor darling.

``I am now well assured that she was.

``We were now in one of the salons. My poor dear child had been dancing,
and was resting a little in one of the chairs near the door; I was
standing near. The two ladies I have mentioned had approached and the
younger took the chair next my ward; while her companion stood beside
me, and for a little time addressed herself, in a low tone, to
her charge.

``Availing herself of the privilege of her mask, she turned to me, and in
the tone of an old friend, and calling me by my name, opened a
conversation with me, which piqued my curiosity a good deal. She
referred to many scenes where she had met me--at Court, and at
distinguished houses. She alluded to little incidents which I had long
ceased to think of, but which, I found, had only lain in abeyance in my
memory, for they instantly started into life at her touch.

``I became more and more curious to ascertain who she was, every moment.
She parried my attempts to discover very adroitly and pleasantly. The
knowledge she showed of many passages in my life seemed to me all but
unaccountable; and she appeared to take a not unnatural pleasure in
foiling my curiosity, and in seeing me flounder in my eager perplexity,
from one conjecture to another.

``In the meantime the young lady, whom her mother called by the odd name
of Millarca, when she once or twice addressed her, had, with the same
ease and grace, got into conversation with my ward.

``She introduced herself by saying that her mother was a very old
acquaintance of mine. She spoke of the agreeable audacity which a mask
rendered practicable; she talked like a friend; she admired her dress,
and insinuated very prettily her admiration of her beauty. She amused
her with laughing criticisms upon the people who crowded the ballroom,
and laughed at my poor child's fun. She was very witty and lively when
she pleased, and after a time they had grown very good friends, and the
young stranger lowered her mask, displaying a remarkably beautiful face.
I had never seen it before, neither had my dear child. But though it was
new to us, the features were so engaging, as well as lovely, that it
was impossible not to feel the attraction powerfully. My poor girl did
so. I never saw anyone more taken with another at first sight, unless,
indeed, it was the stranger herself, who seemed quite to have lost her
heart to her.

``In the meantime, availing myself of the license of a masquerade, I put
not a few questions to the elder lady.

``{}`You have puzzled me utterly,' I said, laughing. `Is that not enough?
Won't you, now, consent to stand on equal terms, and do me the kindness
to remove your mask?'

``{}`Can any request be more unreasonable?' she replied. `Ask a lady to
yield an advantage! Beside, how do you know you should recognize me?
Years make changes.'

``{}`As you see,' I said, with a bow, and, I suppose, a rather melancholy
little laugh.

``{}`As philosophers tell us,' she said; `and how do you know that a sight
of my face would help you?'

``{}`I should take chance for that,' I answered. `It is vain trying to make
yourself out an old woman; your figure betrays you.'

``{}`Years, nevertheless, have passed since I saw you, rather since you saw
me, for that is what I am considering. Millarca, there, is my daughter;
I cannot then be young, even in the opinion of people whom time has
taught to be indulgent, and I may not like to be compared with what you
remember me. You have no mask to remove. You can offer me nothing in
exchange.'

``{}`My petition is to your pity, to remove it.'

``{}`And mine to yours, to let it stay where it is,' she replied.

``{}`Well, then, at least you will tell me whether you are French or
German; you speak both languages so perfectly.'

``{}`I don't think I shall tell you that, General; you intend a surprise,
and are meditating the particular point of attack.'

``{}`At all events, you won't deny this,' I said, `that being honored by
your permission to converse, I ought to know how to address you. Shall I
say Madame la Comtesse?'

``She laughed, and she would, no doubt, have met me with another
evasion--if, indeed, I can treat any occurrence in an interview every
circumstance of which was prearranged, as I now believe, with the
profoundest cunning, as liable to be modified by accident.

``{}`As to that,' she began; but she was interrupted, almost as she opened
her lips, by a gentleman, dressed in black, who looked particularly
elegant and distinguished, with this drawback, that his face was the
most deadly pale I ever saw, except in death. He was in no
masquerade--in the plain evening dress of a gentleman; and he said,
without a smile, but with a courtly and unusually low bow:--

``{}`Will Madame la Comtesse permit me to say a very few words which may
interest her?'

``The lady turned quickly to him, and touched her lip in token of
silence; she then said to me, `Keep my place for me, General; I shall
return when I have said a few words.'

``And with this injunction, playfully given, she walked a little aside
with the gentleman in black, and talked for some minutes, apparently
very earnestly. They then walked away slowly together in the crowd, and
I lost them for some minutes.

``I spent the interval in cudgeling my brains for a conjecture as to the
identity of the lady who seemed to remember me so kindly, and I was
thinking of turning about and joining in the conversation between my
pretty ward and the Countess's daughter, and trying whether, by the time
she returned, I might not have a surprise in store for her, by having
her name, title, chateau, and estates at my fingers' ends. But at this
moment she returned, accompanied by the pale man in black, who said:

``{}`I shall return and inform Madame la Comtesse when her carriage is at
the door.'

``He withdrew with a bow.''

\pbreak{}

\chapter{A Petition}\hypertarget{a-petition}{}\label{a-petition}

``{}`Then we are to lose Madame la Comtesse, but I hope only for a few
hours,' I said, with a low bow.

``{}`It may be that only, or it may be a few weeks. It was very unlucky his
speaking to me just now as he did. Do you now know me?'

``I assured her I did not.

``{}`You shall know me,' she said, `but not at present. We are older and
better friends than, perhaps, you suspect. I cannot yet declare myself.
I shall in three weeks pass your beautiful schloss, about which I have
been making enquiries. I shall then look in upon you for an hour or two,
and renew a friendship which I never think of without a thousand
pleasant recollections. This moment a piece of news has reached me like
a thunderbolt. I must set out now, and travel by a devious route, nearly
a hundred miles, with all the dispatch I can possibly make. My
perplexities multiply. I am only deterred by the compulsory reserve I
practice as to my name from making a very singular request of you. My
poor child has not quite recovered her strength. Her horse fell with
her, at a hunt which she had ridden out to witness, her nerves have not
yet recovered the shock, and our physician says that she must on no
account exert herself for some time to come. We came here, in
consequence, by very easy stages--hardly six leagues a day. I must now
travel day and night, on a mission of life and death--a mission the
critical and momentous nature of which I shall be able to explain to you
when we meet, as I hope we shall, in a few weeks, without the necessity
of any concealment.'

``She went on to make her petition, and it was in the tone of a person
from whom such a request amounted to conferring, rather than seeking
a favor.

``This was only in manner, and, as it seemed, quite unconsciously. Than
the terms in which it was expressed, nothing could be more deprecatory.
It was simply that I would consent to take charge of her daughter during
her absence.

``This was, all things considered, a strange, not to say, an audacious
request. She in some sort disarmed me, by stating and admitting
everything that could be urged against it, and throwing herself entirely
upon my chivalry. At the same moment, by a fatality that seems to have
predetermined all that happened, my poor child came to my side, and, in
an undertone, besought me to invite her new friend, Millarca, to pay us
a visit. She had just been sounding her, and thought, if her mamma would
allow her, she would like it extremely.

``At another time I should have told her to wait a little, until, at
least, we knew who they were. But I had not a moment to think in. The
two ladies assailed me together, and I must confess the refined and
beautiful face of the young lady, about which there was something
extremely engaging, as well as the elegance and fire of high birth,
determined me; and, quite overpowered, I submitted, and undertook, too
easily, the care of the young lady, whom her mother called Millarca.

``The Countess beckoned to her daughter, who listened with grave
attention while she told her, in general terms, how suddenly and
peremptorily she had been summoned, and also of the arrangement she had
made for her under my care, adding that I was one of her earliest and
most valued friends.

``I made, of course, such speeches as the case seemed to call for, and
found myself, on reflection, in a position which I did not half like.

``The gentleman in black returned, and very ceremoniously conducted the
lady from the room.

``The demeanor of this gentleman was such as to impress me with the
conviction that the Countess was a lady of very much more importance
than her modest title alone might have led me to assume.

``Her last charge to me was that no attempt was to be made to learn more
about her than I might have already guessed, until her return. Our
distinguished host, whose guest she was, knew her reasons.

``{}`But here,' she said, `neither I nor my daughter could safely remain
for more than a day. I removed my mask imprudently for a moment, about
an hour ago, and, too late, I fancied you saw me. So I resolved to seek
an opportunity of talking a little to you. Had I found that you had seen
me, I would have thrown myself on your high sense of honor to keep my
secret some weeks. As it is, I am satisfied that you did not see me; but
if you now suspect, or, on reflection, should suspect, who I am, I
commit myself, in like manner, entirely to your honor. My daughter will
observe the same secrecy, and I well know that you will, from time to
time, remind her, lest she should thoughtlessly disclose it.'

``She whispered a few words to her daughter, kissed her hurriedly twice,
and went away, accompanied by the pale gentleman in black, and
disappeared in the crowd.

``{}`In the next room,' said Millarca, `there is a window that looks upon
the hall door. I should like to see the last of mamma, and to kiss my
hand to her.'

``We assented, of course, and accompanied her to the window. We looked
out, and saw a handsome old-fashioned carriage, with a troop of couriers
and footmen. We saw the slim figure of the pale gentleman in black, as
he held a thick velvet cloak, and placed it about her shoulders and
threw the hood over her head. She nodded to him, and just touched his
hand with hers. He bowed low repeatedly as the door closed, and the
carriage began to move.

``{}`She is gone,' said Millarca, with a sigh.

``{}`She is gone,' I repeated to myself, for the first time--in the hurried
moments that had elapsed since my consent--reflecting upon the folly
of my act.

``{}`She did not look up,' said the young lady, plaintively.

``{}`The Countess had taken off her mask, perhaps, and did not care to show
her face,' I said; `and she could not know that you were in the window.'

``She sighed, and looked in my face. She was so beautiful that I
relented. I was sorry I had for a moment repented of my hospitality, and
I determined to make her amends for the unavowed churlishness of my
reception.

``The young lady, replacing her mask, joined my ward in persuading me to
return to the grounds, where the concert was soon to be renewed. We did
so, and walked up and down the terrace that lies under the
castle windows.

``Millarca became very intimate with us, and amused us with lively
descriptions and stories of most of the great people whom we saw upon
the terrace. I liked her more and more every minute. Her gossip without
being ill-natured, was extremely diverting to me, who had been so long
out of the great world. I thought what life she would give to our
sometimes lonely evenings at home.

``This ball was not over until the morning sun had almost reached the
horizon. It pleased the Grand Duke to dance till then, so loyal people
could not go away, or think of bed.

``We had just got through a crowded saloon, when my ward asked me what
had become of Millarca. I thought she had been by her side, and she
fancied she was by mine. The fact was, we had lost her.

``All my efforts to find her were vain. I feared that she had mistaken,
in the confusion of a momentary separation from us, other people for her
new friends, and had, possibly, pursued and lost them in the extensive
grounds which were thrown open to us.

``Now, in its full force, I recognized a new folly in my having
undertaken the charge of a young lady without so much as knowing her
name; and fettered as I was by promises, of the reasons for imposing
which I knew nothing, I could not even point my inquiries by saying that
the missing young lady was the daughter of the Countess who had taken
her departure a few hours before.

``Morning broke. It was clear daylight before I gave up my search. It was
not till near two o'clock next day that we heard anything of my
missing charge.

``At about that time a servant knocked at my niece's door, to say that he
had been earnestly requested by a young lady, who appeared to be in
great distress, to make out where she could find the General Baron
Spielsdorf and the young lady his daughter, in whose charge she had been
left by her mother.

``There could be no doubt, notwithstanding the slight inaccuracy, that
our young friend had turned up; and so she had. Would to heaven we
had lost her!

``She told my poor child a story to account for her having failed to
recover us for so long. Very late, she said, she had got to the
housekeeper's bedroom in despair of finding us, and had then fallen
into a deep sleep which, long as it was, had hardly sufficed to recruit
her strength after the fatigues of the ball.

``That day Millarca came home with us. I was only too happy, after all,
to have secured so charming a companion for my dear girl.''

\pbreak{}

\chapter{The Woodman}\hypertarget{the-woodman}{}\label{the-woodman}

``There soon, however, appeared some drawbacks. In the first place,
Millarca complained of extreme languor--the weakness that remained after
her late illness--and she never emerged from her room till the afternoon
was pretty far advanced. In the next place, it was accidentally
discovered, although she always locked her door on the inside, and never
disturbed the key from its place till she admitted the maid to assist at
her toilet, that she was undoubtedly sometimes absent from her room in
the very early morning, and at various times later in the day, before
she wished it to be understood that she was stirring. She was repeatedly
seen from the windows of the schloss, in the first faint grey of the
morning, walking through the trees, in an easterly direction, and
looking like a person in a trance. This convinced me that she walked in
her sleep. But this hypothesis did not solve the puzzle. How did she
pass out from her room, leaving the door locked on the inside? How did
she escape from the house without unbarring door or window?

``In the midst of my perplexities, an anxiety of a far more urgent kind
presented itself.

``My dear child began to lose her looks and health, and that in a manner
so mysterious, and even horrible, that I became thoroughly frightened.

``She was at first visited by appalling dreams; then, as she fancied, by
a specter, sometimes resembling Millarca, sometimes in the shape of a
beast, indistinctly seen, walking round the foot of her bed, from
side to side.

``Lastly came sensations. One, not unpleasant, but very peculiar, she
said, resembled the flow of an icy stream against her breast. At a later
time, she felt something like a pair of large needles pierce her, a
little below the throat, with a very sharp pain. A few nights after,
followed a gradual and convulsive sense of strangulation; then came
unconsciousness.''

I could hear distinctly every word the kind old General was saying,
because by this time we were driving upon the short grass that spreads
on either side of the road as you approach the roofless village which
had not shown the smoke of a chimney for more than half a century.

You may guess how strangely I felt as I heard my own symptoms so exactly
described in those which had been experienced by the poor girl who, but
for the catastrophe which followed, would have been at that moment a
visitor at my father's chateau. You may suppose, also, how I felt as I
heard him detail habits and mysterious peculiarities which were, in
fact, those of our beautiful guest, Carmilla!

A vista opened in the forest; we were on a sudden under the chimneys and
gables of the ruined village, and the towers and battlements of the
dismantled castle, round which gigantic trees are grouped, overhung us
from a slight eminence.

In a frightened dream I got down from the carriage, and in silence, for
we had each abundant matter for thinking; we soon mounted the ascent,
and were among the spacious chambers, winding stairs, and dark
corridors of the castle.

``And this was once the palatial residence of the Karnsteins!'' said the
old General at length, as from a great window he looked out across the
village, and saw the wide, undulating expanse of forest. ``It was a bad
family, and here its bloodstained annals were written,'' he continued.
``It is hard that they should, after death, continue to plague the human
race with their atrocious lusts. That is the chapel of the Karnsteins,
down there.''

He pointed down to the grey walls of the Gothic building partly visible
through the foliage, a little way down the steep. ``And I hear the axe of
a woodman,'' he added, ``busy among the trees that surround it; he
possibly may give us the information of which I am in search, and point
out the grave of Mircalla, Countess of Karnstein. These rustics preserve
the local traditions of great families, whose stories die out among the
rich and titled so soon as the families themselves become extinct.''

``We have a portrait, at home, of Mircalla, the Countess Karnstein;
should you like to see it?'' asked my father.

``Time enough, dear friend,'' replied the General. ``I believe that I have
seen the original; and one motive which has led me to you earlier than I
at first intended, was to explore the chapel which we are now
approaching.''

``What! see the Countess Mircalla,'' exclaimed my father; ``why, she has
been dead more than a century!''

``Not so dead as you fancy, I am told,'' answered the General.

``I confess, General, you puzzle me utterly,'' replied my father, looking
at him, I fancied, for a moment with a return of the suspicion I
detected before. But although there was anger and detestation, at times,
in the old General's manner, there was nothing flighty.

``There remains to me,'' he said, as we passed under the heavy arch of
the Gothic church--for its dimensions would have justified its being so
styled--``but one object which can interest me during the few years that
remain to me on earth, and that is to wreak on her the vengeance which,
I thank God, may still be accomplished by a mortal arm.''

``What vengeance can you mean?'' asked my father, in increasing amazement.

``I mean, to decapitate the monster,'' he answered, with a fierce flush,
and a stamp that echoed mournfully through the hollow ruin, and his
clenched hand was at the same moment raised, as if it grasped the handle
of an axe, while he shook it ferociously in the air.

``What?'' exclaimed my father, more than ever bewildered.

``To strike her head off.''

``Cut her head off!''

``Aye, with a hatchet, with a spade, or with anything that can cleave
through her murderous throat. You shall hear,'' he answered, trembling
with rage. And hurrying forward he said:

``That beam will answer for a seat; your dear child is fatigued; let her
be seated, and I will, in a few sentences, close my dreadful story.''

The squared block of wood, which lay on the grass-grown pavement of the
chapel, formed a bench on which I was very glad to seat myself, and in
the meantime the General called to the woodman, who had been removing
some boughs which leaned upon the old walls; and, axe in hand, the hardy
old fellow stood before us.

He could not tell us anything of these monuments; but there was an old
man, he said, a ranger of this forest, at present sojourning in the
house of the priest, about two miles away, who could point out every
monument of the old Karnstein family; and, for a trifle, he undertook
to bring him back with him, if we would lend him one of our horses, in
little more than half an hour.

``Have you been long employed about this forest?'' asked my father of the
old man.

``I have been a woodman here,'' he answered in his patois, ``under the
forester, all my days; so has my father before me, and so on, as many
generations as I can count up. I could show you the very house in the
village here, in which my ancestors lived.''

``How came the village to be deserted?'' asked the General.

``It was troubled by revenants, sir; several were tracked to their
graves, there detected by the usual tests, and extinguished in the usual
way, by decapitation, by the stake, and by burning; but not until many
of the villagers were killed.

``But after all these proceedings according to law,'' he continued--``so
many graves opened, and so many vampires deprived of their horrible
animation--the village was not relieved. But a Moravian nobleman, who
happened to be traveling this way, heard how matters were, and being
skilled--as many people are in his country--in such affairs, he offered
to deliver the village from its tormentor. He did so thus: There being a
bright moon that night, he ascended, shortly after sunset, the towers of
the chapel here, from whence he could distinctly see the churchyard
beneath him; you can see it from that window. From this point he watched
until he saw the vampire come out of his grave, and place near it the
linen clothes in which he had been folded, and then glide away towards
the village to plague its inhabitants.

``The stranger, having seen all this, came down from the steeple, took
the linen wrappings of the vampire, and carried them up to the top of
the tower, which he again mounted. When the vampire returned from his
prowlings and missed his clothes, he cried furiously to the Moravian,
whom he saw at the summit of the tower, and who, in reply, beckoned him
to ascend and take them. Whereupon the vampire, accepting his
invitation, began to climb the steeple, and so soon as he had reached
the battlements, the Moravian, with a stroke of his sword, clove his
skull in twain, hurling him down to the churchyard, whither, descending
by the winding stairs, the stranger followed and cut his head off, and
next day delivered it and the body to the villagers, who duly impaled
and burnt them.

``This Moravian nobleman had authority from the then head of the family
to remove the tomb of Mircalla, Countess Karnstein, which he did
effectually, so that in a little while its site was quite forgotten.''

``Can you point out where it stood?'' asked the General, eagerly.

The forester shook his head, and smiled.

``Not a soul living could tell you that now,'' he said; ``besides, they say
her body was removed; but no one is sure of that either.''

Having thus spoken, as time pressed, he dropped his axe and departed,
leaving us to hear the remainder of the General's strange story.

\pbreak{}

\chapter{The Meeting}\hypertarget{the-meeting}{}\label{the-meeting}

``My beloved child,'' he resumed, ``was now growing rapidly worse. The
physician who attended her had failed to produce the slightest
impression on her disease, for such I then supposed it to be. He saw my
alarm, and suggested a consultation. I called in an abler physician,
from Gratz.

``Several days elapsed before he arrived. He was a good and pious, as well
as a learned man. Having seen my poor ward together, they withdrew to my
library to confer and discuss. I, from the adjoining room, where I
awaited their summons, heard these two gentlemen's voices raised in
something sharper than a strictly philosophical discussion. I knocked at
the door and entered. I found the old physician from Gratz maintaining
his theory. His rival was combating it with undisguised ridicule,
accompanied with bursts of laughter. This unseemly manifestation
subsided and the altercation ended on my entrance.

``{}`Sir,' said my first physician, `my learned brother seems to think that
you want a conjuror, and not a doctor.'

``{}`Pardon me,' said the old physician from Gratz, looking displeased, `I
shall state my own view of the case in my own way another time. I
grieve, Monsieur le General, that by my skill and science I can be of no
use. Before I go I shall do myself the honor to suggest something to
you.'

``He seemed thoughtful, and sat down at a table and began to write.

``Profoundly disappointed, I made my bow, and as I turned to go, the other
doctor pointed over his shoulder to his companion who was writing, and
then, with a shrug, significantly touched his forehead.

``This consultation, then, left me precisely where I was. I walked out
into the grounds, all but distracted. The doctor from Gratz, in ten or
fifteen minutes, overtook me. He apologized for having followed me, but
said that he could not conscientiously take his leave without a few
words more. He told me that he could not be mistaken; no natural disease
exhibited the same symptoms; and that death was already very near. There
remained, however, a day, or possibly two, of life. If the fatal seizure
were at once arrested, with great care and skill her strength might
possibly return. But all hung now upon the confines of the irrevocable.
One more assault might extinguish the last spark of vitality which is,
every moment, ready to die.

``{}`And what is the nature of the seizure you speak of?' I entreated.

``{}`I have stated all fully in this note, which I place in your hands upon
the distinct condition that you send for the nearest clergyman, and open
my letter in his presence, and on no account read it till he is with
you; you would despise it else, and it is a matter of life and death.
Should the priest fail you, then, indeed, you may read it.'

``He asked me, before taking his leave finally, whether I would wish to
see a man curiously learned upon the very subject, which, after I had
read his letter, would probably interest me above all others, and he
urged me earnestly to invite him to visit him there; and so took
his leave.

``The ecclesiastic was absent, and I read the letter by myself. At
another time, or in another case, it might have excited my ridicule. But
into what quackeries will not people rush for a last chance, where all
accustomed means have failed, and the life of a beloved object is
at stake?

``Nothing, you will say, could be more absurd than the learned man's
letter.

``It was monstrous enough to have consigned him to a madhouse. He said
that the patient was suffering from the visits of a vampire! The
punctures which she described as having occurred near the throat, were,
he insisted, the insertion of those two long, thin, and sharp teeth
which, it is well known, are peculiar to vampires; and there could be no
doubt, he added, as to the well-defined presence of the small livid mark
which all concurred in describing as that induced by the demon's lips,
and every symptom described by the sufferer was in exact conformity with
those recorded in every case of a similar visitation.

``Being myself wholly skeptical as to the existence of any such portent
as the vampire, the supernatural theory of the good doctor furnished, in
my opinion, but another instance of learning and intelligence oddly
associated with some one hallucination. I was so miserable, however,
that, rather than try nothing, I acted upon the instructions of
the letter.

``I concealed myself in the dark dressing room, that opened upon the poor
patient's room, in which a candle was burning, and watched there till
she was fast asleep. I stood at the door, peeping through the small
crevice, my sword laid on the table beside me, as my directions
prescribed, until, a little after one, I saw a large black object, very
ill-defined, crawl, as it seemed to me, over the foot of the bed, and
swiftly spread itself up to the poor girl's throat, where it swelled, in
a moment, into a great, palpitating mass.

``For a few moments I had stood petrified. I now sprang forward, with my
sword in my hand. The black creature suddenly contracted towards the
foot of the bed, glided over it, and, standing on the floor about a yard
below the foot of the bed, with a glare of skulking ferocity and horror
fixed on me, I saw Millarca. Speculating I know not what, I struck at
her instantly with my sword; but I saw her standing near the door,
unscathed. Horrified, I pursued, and struck again. She was gone; and my
sword flew to shivers against the door.

``I can't describe to you all that passed on that horrible night. The
whole house was up and stirring. The specter Millarca was gone. But her
victim was sinking fast, and before the morning dawned, she died.''

The old General was agitated. We did not speak to him. My father walked
to some little distance, and began reading the inscriptions on the
tombstones; and thus occupied, he strolled into the door of a side
chapel to prosecute his researches. The General leaned against the wall,
dried his eyes, and sighed heavily. I was relieved on hearing the voices
of Carmilla and Madame, who were at that moment approaching. The voices
died away.

In this solitude, having just listened to so strange a story, connected,
as it was, with the great and titled dead, whose monuments were
moldering among the dust and ivy round us, and every incident of which
bore so awfully upon my own mysterious case--in this haunted spot,
darkened by the towering foliage that rose on every side, dense and high
above its noiseless walls--a horror began to steal over me, and my heart
sank as I thought that my friends were, after all, not about to enter
and disturb this triste and ominous scene.

The old General's eyes were fixed on the ground, as he leaned with his
hand upon the basement of a shattered monument.

Under a narrow, arched doorway, surmounted by one of those demoniacal
grotesques in which the cynical and ghastly fancy of old Gothic carving
delights, I saw very gladly the beautiful face and figure of Carmilla
enter the shadowy chapel.

I was just about to rise and speak, and nodded smiling, in answer to her
peculiarly engaging smile; when with a cry, the old man by my side
caught up the woodman's hatchet, and started forward. On seeing him a
brutalized change came over her features. It was an instantaneous and
horrible transformation, as she made a crouching step backwards. Before
I could utter a scream, he struck at her with all his force, but she
dived under his blow, and unscathed, caught him in her tiny grasp by the
wrist. He struggled for a moment to release his arm, but his hand
opened, the axe fell to the ground, and the girl was gone.

He staggered against the wall. His grey hair stood upon his head, and a
moisture shone over his face, as if he were at the point of death.

The frightful scene had passed in a moment. The first thing I recollect
after, is Madame standing before me, and impatiently repeating again and
again, the question, ``Where is Mademoiselle Carmilla?''

I answered at length, ``I don't know--I can't tell--she went there,'' and
I pointed to the door through which Madame had just entered; ``only a
minute or two since.''

``But I have been standing there, in the passage, ever since Mademoiselle
Carmilla entered; and she did not return.''

She then began to call ``Carmilla,'' through every door and passage and
from the windows, but no answer came.

``She called herself Carmilla?'' asked the General, still agitated.

``Carmilla, yes,'' I answered.

``Aye,'' he said; ``that is Millarca. That is the same person who long ago
was called Mircalla, Countess Karnstein. Depart from this accursed
ground, my poor child, as quickly as you can. Drive to the clergyman's
house, and stay there till we come. Begone! May you never behold
Carmilla more; you will not find her here.''

\pbreak{}

\chapter{Ordeal and Execution}\hypertarget{ordeal-and-execution}{}\label{ordeal-and-execution}

As he spoke one of the strangest looking men I ever beheld entered the
chapel at the door through which Carmilla had made her entrance and her
exit. He was tall, narrow-chested, stooping, with high shoulders, and
dressed in black. His face was brown and dried in with deep furrows; he
wore an oddly-shaped hat with a broad leaf. His hair, long and grizzled,
hung on his shoulders. He wore a pair of gold spectacles, and walked
slowly, with an odd shambling gait, with his face sometimes turned up to
the sky, and sometimes bowed down towards the ground, seemed to wear a
perpetual smile; his long thin arms were swinging, and his lank hands,
in old black gloves ever so much too wide for them, waving and
gesticulating in utter abstraction.

``The very man!'' exclaimed the General, advancing with manifest delight.
``My dear Baron, how happy I am to see you, I had no hope of meeting you
so soon.'' He signed to my father, who had by this time returned, and
leading the fantastic old gentleman, whom he called the Baron to meet
him. He introduced him formally, and they at once entered into earnest
conversation. The stranger took a roll of paper from his pocket, and
spread it on the worn surface of a tomb that stood by. He had a pencil
case in his fingers, with which he traced imaginary lines from point to
point on the paper, which from their often glancing from it, together,
at certain points of the building, I concluded to be a plan of the
chapel. He accompanied, what I may term, his lecture, with occasional
readings from a dirty little book, whose yellow leaves were closely
written over.

They sauntered together down the side aisle, opposite to the spot where
I was standing, conversing as they went; then they began measuring
distances by paces, and finally they all stood together, facing a piece
of the sidewall, which they began to examine with great minuteness;
pulling off the ivy that clung over it, and rapping the plaster with the
ends of their sticks, scraping here, and knocking there. At length they
ascertained the existence of a broad marble tablet, with letters carved
in relief upon it.

With the assistance of the woodman, who soon returned, a monumental
inscription, and carved escutcheon, were disclosed. They proved to be
those of the long lost monument of Mircalla, Countess Karnstein.

The old General, though not I fear given to the praying mood, raised his
hands and eyes to heaven, in mute thanksgiving for some moments.

``Tomorrow,'' I heard him say; ``the commissioner will be here, and the
Inquisition will be held according to law.''

Then turning to the old man with the gold spectacles, whom I have
described, he shook him warmly by both hands and said:

``Baron, how can I thank you? How can we all thank you? You will have
delivered this region from a plague that has scourged its inhabitants
for more than a century. The horrible enemy, thank God, is at
last tracked.''

My father led the stranger aside, and the General followed. I know that
he had led them out of hearing, that he might relate my case, and I saw
them glance often quickly at me, as the discussion proceeded.

My father came to me, kissed me again and again, and leading me from the
chapel, said:

``It is time to return, but before we go home, we must add to our party
the good priest, who lives but a little way from this; and persuade him
to accompany us to the schloss.''

In this quest we were successful: and I was glad, being unspeakably
fatigued when we reached home. But my satisfaction was changed to
dismay, on discovering that there were no tidings of Carmilla. Of the
scene that had occurred in the ruined chapel, no explanation was offered
to me, and it was clear that it was a secret which my father for the
present determined to keep from me.

The sinister absence of Carmilla made the remembrance of the scene more
horrible to me. The arrangements for the night were singular. Two
servants, and Madame were to sit up in my room that night; and the
ecclesiastic with my father kept watch in the adjoining dressing room.

The priest had performed certain solemn rites that night, the purport of
which I did not understand any more than I comprehended the reason of
this extraordinary precaution taken for my safety during sleep.

I saw all clearly a few days later.

The disappearance of Carmilla was followed by the discontinuance of my
nightly sufferings.

You have heard, no doubt, of the appalling superstition that prevails in
Upper and Lower Styria, in Moravia, Silesia, in Turkish Serbia, in
Poland, even in Russia; the superstition, so we must call it, of
the Vampire.

If human testimony, taken with every care and solemnity, judicially,
before commissions innumerable, each consisting of many members, all
chosen for integrity and intelligence, and constituting reports more
voluminous perhaps than exist upon any one other class of cases, is
worth anything, it is difficult to deny, or even to doubt the existence
of such a phenomenon as the Vampire.

For my part I have heard no theory by which to explain what I myself
have witnessed and experienced, other than that supplied by the ancient
and well-attested belief of the country.

The next day the formal proceedings took place in the Chapel of
Karnstein.

The grave of the Countess Mircalla was opened; and the General and my
father recognized each his perfidious and beautiful guest, in the face
now disclosed to view. The features, though a hundred and fifty years
had passed since her funeral, were tinted with the warmth of life. Her
eyes were open; no cadaverous smell exhaled from the coffin. The two
medical men, one officially present, the other on the part of the
promoter of the inquiry, attested the marvelous fact that there was a
faint but appreciable respiration, and a corresponding action of the
heart. The limbs were perfectly flexible, the flesh elastic; and the
leaden coffin floated with blood, in which to a depth of seven inches,
the body lay immersed.

Here then, were all the admitted signs and proofs of vampirism. The
body, therefore, in accordance with the ancient practice, was raised,
and a sharp stake driven through the heart of the vampire, who uttered a
piercing shriek at the moment, in all respects such as might escape from
a living person in the last agony. Then the head was struck off, and a
torrent of blood flowed from the severed neck. The body and head was
next placed on a pile of wood, and reduced to ashes, which were thrown
upon the river and borne away, and that territory has never since been
plagued by the visits of a vampire.

My father has a copy of the report of the Imperial Commission, with the
signatures of all who were present at these proceedings, attached in
verification of the statement. It is from this official paper that I
have summarized my account of this last shocking scene.

\pbreak{}

\chapter{Conclusion}\hypertarget{conclusion}{}\label{conclusion}

I write all this you suppose with composure. But far from it; I cannot
think of it without agitation. Nothing but your earnest desire so
repeatedly expressed, could have induced me to sit down to a task that
has unstrung my nerves for months to come, and reinduced a shadow of the
unspeakable horror which years after my deliverance continued to make my
days and nights dreadful, and solitude insupportably terrific.

Let me add a word or two about that quaint Baron Vordenburg, to whose
curious lore we were indebted for the discovery of the Countess
Mircalla's grave.

He had taken up his abode in Gratz, where, living upon a mere pittance,
which was all that remained to him of the once princely estates of his
family, in Upper Styria, he devoted himself to the minute and laborious
investigation of the marvelously authenticated tradition of Vampirism.
He had at his fingers' ends all the great and little works upon
the subject.

``Magia Posthuma,'' ``Phlegon de Mirabilibus,'' ``Augustinus de cura pro
Mortuis,'' ``Philosophicae et Christianae Cogitationes de Vampiris,'' by
John Christofer Herenberg; and a thousand others, among which I
remember only a few of those which he lent to my father. He had a
voluminous digest of all the judicial cases, from which he had extracted
a system of principles that appear to govern--some always, and others
occasionally only--the condition of the vampire. I may mention, in
passing, that the deadly pallor attributed to that sort of revenants, is
a mere melodramatic fiction. They present, in the grave, and when they
show themselves in human society, the appearance of healthy life. When
disclosed to light in their coffins, they exhibit all the symptoms that
are enumerated as those which proved the vampire-life of the long-dead
Countess Karnstein.

How they escape from their graves and return to them for certain hours
every day, without displacing the clay or leaving any trace of
disturbance in the state of the coffin or the cerements, has always been
admitted to be utterly inexplicable. The amphibious existence of the
vampire is sustained by daily renewed slumber in the grave. Its horrible
lust for living blood supplies the vigor of its waking existence. The
vampire is prone to be fascinated with an engrossing vehemence,
resembling the passion of love, by particular persons. In pursuit of
these it will exercise inexhaustible patience and stratagem, for access
to a particular object may be obstructed in a hundred ways. It will
never desist until it has satiated its passion, and drained the very
life of its coveted victim. But it will, in these cases, husband and
protract its murderous enjoyment with the refinement of an epicure, and
heighten it by the gradual approaches of an artful courtship. In these
cases it seems to yearn for something like sympathy and consent. In
ordinary ones it goes direct to its object, overpowers with violence,
and strangles and exhausts often at a single feast.

The vampire is, apparently, subject, in certain situations, to special
conditions. In the particular instance of which I have given you a
relation, Mircalla seemed to be limited to a name which, if not her real
one, should at least reproduce, without the omission or addition of a
single letter, those, as we say, anagrammatically, which compose it.

Carmilla did this; so did Millarca.

My father related to the Baron Vordenburg, who remained with us for two
or three weeks after the expulsion of Carmilla, the story about the
Moravian nobleman and the vampire at Karnstein churchyard, and then he
asked the Baron how he had discovered the exact position of the
long-concealed tomb of the Countess Mircalla? The Baron's grotesque
features puckered up into a mysterious smile; he looked down, still
smiling on his worn spectacle case and fumbled with it. Then looking
up, he said:

``I have many journals, and other papers, written by that remarkable man;
the most curious among them is one treating of the visit of which you
speak, to Karnstein. The tradition, of course, discolors and distorts a
little. He might have been termed a Moravian nobleman, for he had
changed his abode to that territory, and was, beside, a noble. But he
was, in truth, a native of Upper Styria. It is enough to say that in
very early youth he had been a passionate and favored lover of the
beautiful Mircalla, Countess Karnstein. Her early death plunged him into
inconsolable grief. It is the nature of vampires to increase and
multiply, but according to an ascertained and ghostly law.

``Assume, at starting, a territory perfectly free from that pest. How
does it begin, and how does it multiply itself? I will tell you. A
person, more or less wicked, puts an end to himself. A suicide, under
certain circumstances, becomes a vampire. That specter visits living
people in their slumbers; they die, and almost invariably, in the grave,
develop into vampires. This happened in the case of the beautiful
Mircalla, who was haunted by one of those demons. My ancestor,
Vordenburg, whose title I still bear, soon discovered this, and in the
course of the studies to which he devoted himself, learned a great
deal more.

``Among other things, he concluded that suspicion of vampirism would
probably fall, sooner or later, upon the dead Countess, who in life had
been his idol. He conceived a horror, be she what she might, of her
remains being profaned by the outrage of a posthumous execution. He has
left a curious paper to prove that the vampire, on its expulsion from
its amphibious existence, is projected into a far more horrible life;
and he resolved to save his once beloved Mircalla from this.

``He adopted the stratagem of a journey here, a pretended removal of her
remains, and a real obliteration of her monument. When age had stolen
upon him, and from the vale of years, he looked back on the scenes he
was leaving, he considered, in a different spirit, what he had done, and
a horror took possession of him. He made the tracings and notes which
have guided me to the very spot, and drew up a confession of the
deception that he had practiced. If he had intended any further action
in this matter, death prevented him; and the hand of a remote descendant
has, too late for many, directed the pursuit to the lair of the beast.''

We talked a little more, and among other things he said was this:

``One sign of the vampire is the power of the hand. The slender hand of
Mircalla closed like a vice of steel on the General's wrist when he
raised the hatchet to strike. But its power is not confined to its
grasp; it leaves a numbness in the limb it seizes, which is slowly, if
ever, recovered from.''

The following Spring my father took me a tour through Italy. We remained
away for more than a year. It was long before the terror of recent
events subsided; and to this hour the image of Carmilla returns to
memory with ambiguous alternations--sometimes the playful, languid,
beautiful girl; sometimes the writhing fiend I saw in the ruined church;
and often from a reverie I have started, fancying I heard the light step
of Carmilla at the drawing room door.

\pbreak{}



% Back Matter
%-------------
\backmatter

%% Bibliography
% \bibliographystyle{\mybibliostyle}
% \bibliocommand
% \pagestyle{empty}
\begin{center}
  \vspace*{\fill}

  {\Large \par
  \noindent Thank you for reading

  \vspace*{0.125in}
  {\LARGE\textbf{\textsf{\thetitle}}}
  \vspace*{0.125in}

  \ifdefined\ReturningCharacters
    \par
    \noindent\ReturningCharacters{} will return in \emph{\NextBook}.} \vspace*{\fill}
  \fi

  {\LARGE\textbf{\textsf{Follow Me}}}

  \vspace*{0.125in}

  \textbf{\theauthor{}} can be followed on Twitter and Facebook.
  \newline Visit his web site at http://dausha.net for more information. You can also sign up for announcements of other books in this series at http://dausha.net.

  \vspace*{\fill}
\end{center}



\end{document}